\documentclass[UTF8, 12pt, fontset=fandol, oneside]{ctexart}
\usepackage{researchnotes}

\title{第五章\quad 多项式}
\author{}
\date{}

\begin{document}

\maketitle

\section*{§5.1\quad 一元多项式代数}

在这一章里，我们将暂时中止对线性空间、线性变换的讨论转而研究多项式. 多项式理论不仅对进一步研究线性代数是必要的，而且在数学的其他分支领域也有极其重要的应用.

\noindent\textbf{定义 5.1.1}\quad 设 $\mathbb{K}$ 是数域，$x$ 为一个形式符号 (称为未定元)，若 $a_0,a_1,\cdots,a_n\in\mathbb{K}(a_n\neq 0,n\geq 0)$，称形式表达式
\[
a_nx^n+a_{n-1}x^{n-1}+\cdots+a_1x+a_0
\]
为数域 $\mathbb{K}$ 上关于未定元 $x$ 的一元 $n$ 次多项式. $\mathbb{K}$ 上的一元多项式全体记为 $\mathbb{K}[x]$.

多项式通常记为 $f(x),g(x)$ 等，其中 $a_ix^i$ 称为第 $i$ 次项，$a_i$ 称为第 $i$ 次项的系数. 如果多项式 $f(x)=a_nx^n+a_{n-1}x^{n-1}+\cdots+a_1x+a_0$ 中 $a_n\neq 0$，则称 $f(x)$ 是一个 $n$ 次多项式，$a_nx^n$ 为 $f(x)$ 的首项或最高次项，$a_0$ 为常数项. 若 $f(x)=a$，则称 $f(x)$ 为常数多项式，当 $a\neq 0$ 时，称之为零次多项式；当 $a=0$ 时，称之为零多项式，规定其次数为 $-\infty$. 两个多项式相等当且仅当它们次数相同且各次项的系数相等，即若
\begin{equation}
f(x)=a_nx^n+a_{n-1}x^{n-1}+\cdots+a_1x+a_0,
\tag{5.1.1}
\end{equation}
\begin{equation}
g(x)=b_mx^m+b_{m-1}x^{m-1}+\cdots+b_1x+b_0,
\tag{5.1.2}
\end{equation}
则 $f(x)=g(x)$ 当且仅当 $m=n,\ a_i=b_i\ (i=0,1,\cdots,n)$.

现在我们定义 $\mathbb{K}[x]$ 上的运算. 设 $f(x),g(x)\in\mathbb{K}[x]$，适当增加几个系数为零的项，可设
\[
f(x)=a_nx^n+a_{n-1}x^{n-1}+\cdots+a_1x+a_0,
\]
\[
g(x)=b_nx^n+b_{n-1}x^{n-1}+\cdots+b_1x+b_0,
\]
定义
\[
f(x)+g(x)=(a_n+b_n)x^n+(a_{n-1}+b_{n-1})x^{n-1}+\cdots+(a_1+b_1)x+(a_0+b_0),
\]
若 $c\in\mathbb{K}$，定义
\[
cf(x)=ca_nx^n+ca_{n-1}x^{n-1}+\cdots+ca_1x+ca_0,
\]
则 $\mathbb{K}[x]$ 在上述定义下成为 $\mathbb{K}$ 上的线性空间，其中零向量为零多项式. 验证工作并不困难，只需逐条验证线性空间的公理即可，请读者自己完成.

再来定义两个多项式的积. 若 $f(x),g(x)$ 如 (5.1.1) 及 (5.1.2) 所示且 $a_n\neq 0,\ b_m\neq 0$，定义 $f(x)\cdot g(x)=h(x)$，其中 $h(x)$ 是一个 $n+m$ 次多项式，若设
\[
h(x)=c_{n+m}x^{n+m}+c_{n+m-1}x^{n+m-1}+\cdots+c_1x+c_0,
\]
则
\[
\begin{aligned}
c_{n+m}&=a_nb_m,\\
c_{n+m-1}&=a_{n-1}b_m+a_nb_{m-1},\\
&\cdots\cdots\cdots\\
c_k&=\sum_{i+j=k}a_ib_j=a_0b_k+a_1b_{k-1}+\cdots+a_kb_0,\\
&\cdots\cdots\cdots\\
c_0&=a_0b_0.
\end{aligned}
\]
不难验证 $\mathbb{K}[x]$ 中元素的乘积适合下列法则：

(1) $f(x)g(x)=g(x)f(x)$；

(2) $(f(x)g(x))h(x)=f(x)(g(x)h(x))$；

(3) $(f(x)+g(x))h(x)=f(x)h(x)+g(x)h(x)$；

(4) $c(f(x)g(x))=(cf(x))g(x)=f(x)(cg(x))$；

(5) 若把 $c$ 看成是常数多项式，则 $c$ 与 $f(x)$ 的作为多项式的积与 $c$ 作为数乘以 $f(x)$ 的积相同.

由上述法则可知，$\mathbb{K}[x]$ 是数域 $\mathbb{K}$ 上的代数，通常称 $\mathbb{K}[x]$ 为数域 $\mathbb{K}$ 上的一元多项式代数. 由于 $\mathbb{K}[x]$ 中的加法及乘法适合环的条件，因此 $\mathbb{K}[x]$ 也称为 $\mathbb{K}$ 上的一元多项式环.

若 $f(x)$ 是 $n$ 次多项式，则记 $f(x)$ 的次数为
\[
\deg f(x)=n.
\]
显然有下列引理.

\noindent\textbf{引理 5.1.1}\quad 若 $f(x),g(x)\in\mathbb{K}[x]$，则
\[
\deg(f(x)g(x))=\deg f(x)+\deg g(x).
\]

注意上述等式对零多项式也适用. 由这个引理我们可证明下面的命题.

\noindent\textbf{命题 5.1.1}\quad 若 $f(x),g(x)\in\mathbb{K}[x]$ 且 $f(x)\neq 0,\ g(x)\neq 0$，则
\[
f(x)g(x)\neq 0.
\]

\noindent\textbf{证明}\quad $\deg(f(x)g(x))=\deg f(x)+\deg g(x)\geq 0+0=0$，因此
\[
f(x)g(x)\neq 0.\quad \square
\]

\noindent\textbf{推论 5.1.1}\quad 若 $f(x)\neq 0$ 且 $f(x)g(x)=f(x)h(x)$，则
\[
g(x)=h(x).
\]

\noindent\textbf{证明}\quad $f(x)(g(x)-h(x))=0$，但 $f(x)\neq 0$，故必须有 $g(x)-h(x)=0.\quad \square$

下列命题也是显然的.

\noindent\textbf{命题 5.1.2}\quad 设 $f(x),g(x)\in\mathbb{K}[x]$，则

(1) $\deg(cf(x))=\deg f(x),\ 0\neq c\in\mathbb{K}$；

(2) $\deg(f(x)+g(x))\leq \max\{\deg f(x),\deg g(x)\}$.

\section*{§5.2\quad 整除}

\noindent\textbf{定义 5.2.1}\quad 设 $f(x),g(x)\in\mathbb{K}[x]$，若存在 $h(x)\in\mathbb{K}[x]$，使
\[
f(x)=g(x)h(x),
\]
则称 $g(x)$ 是 $f(x)$ 的因式，或 $g(x)$ 可以整除 $f(x)$，或 $f(x)$ 可以被 $g(x)$ 整除，记为 $g(x)\mid f(x)$. 否则称 $g(x)$ 不能整除 $f(x)$，或 $f(x)$ 不能被 $g(x)$ 整除.

从整除的定义知道，零多项式的因式可以是任一多项式. 但零多项式不能是任一非零多项式的因式. 整除有下列简单性质.

\noindent\textbf{命题 5.2.1}\quad 设 $f(x),g(x),h(x)\in\mathbb{K}[x],\ 0\neq c\in\mathbb{K}$，则

(1) 若 $f(x)\mid g(x)$，则 $cf(x)\mid g(x)$，因此非零常数多项式 $c$ 是任一非零多项式的因式；

(2) $f(x)\mid f(x)$；

(3) 若 $f(x)\mid g(x),\ g(x)\mid h(x)$，则 $f(x)\mid h(x)$；

(4) 若 $f(x)\mid g(x),\ f(x)\mid h(x)$，则对任意的多项式 $u(x),v(x)$，有
\[
f(x)\mid g(x)u(x)+h(x)v(x);
\]

(5) 设 $f(x)\mid g(x),\ g(x)\mid f(x)$ 且 $f(x),g(x)$ 都是非零多项式，则存在 $\mathbb{K}$ 中非零元 $c$，使
\[
f(x)=cg(x).
\]

\noindent\textbf{证明}\quad (1) 若 $g(x)=f(x)p(x)$，则
\[
g(x)=(cf(x))(c^{-1}p(x)).
\]
此即 $cf(x)\mid g(x)$.

(2) 显然.

(3) 若 $g(x)=f(x)p(x),\ h(x)=g(x)q(x)$，则
\[
h(x)=(f(x)p(x))q(x)=f(x)(p(x)q(x)).
\]

(4) 若 $g(x)=f(x)p(x),\ h(x)=f(x)q(x)$，则
\[
g(x)u(x)+h(x)v(x)=f(x)(p(x)u(x)+q(x)v(x)).
\]

(5) 设 $g(x)=f(x)p(x),\ f(x)=g(x)q(x)$，则
\[
f(x)=f(x)(p(x)q(x)).
\]
由此即得
\[
\deg f(x)=\deg f(x)+\deg(p(x)q(x)),
\]
从而
\[
\deg(p(x)q(x))=0,
\]
于是
\[
\deg p(x)=\deg q(x)=0.
\]
因此 $p(x)$ 及 $q(x)$ 均为非零常数多项式，即 $f(x)$ 和 $g(x)$ 相差一个非零常数. $\square$

\noindent\textbf{注}\quad 适合命题中条件 (5) 的两个多项式 (即可以互相整除的两个多项式) 称为相伴多项式，记为 $f(x)\sim g(x)$.

任意给定两个非零多项式 $f(x)$ 及 $g(x)$，未必有 $f(x)\mid g(x)$ 或 $g(x)\mid f(x)$，但仍可做带余式的除法.

\noindent\textbf{定理 5.2.1}\quad 设 $f(x),g(x)\in\mathbb{K}[x],\ g(x)\neq 0$，则必存在唯一的 $q(x),r(x)\in\mathbb{K}[x]$，使
\begin{equation}
f(x)=g(x)q(x)+r(x),
\tag{5.2.1}
\end{equation}
且 $\deg r(x)<\deg g(x)$.

\noindent\textbf{证明}\quad 若 $\deg f(x)<\deg g(x)$，只需令 $q(x)=0,\ r(x)=f(x)$ 即可.
现设 $\deg f(x)\geq \deg g(x)$，对 $f(x)$ 的次数用数学归纳法. 若 $\deg f(x)=0$，则 $\deg g(x)=0$. 因此可设 $f(x)=a,\ g(x)=b\ (a\neq 0,b\neq 0)$. 这时令 $q(x)=ab^{-1},\ r(x)=0$ 即可. 作为归纳假设，我们设结论对小于 $n$ 次的多项式均成立. 设
\[
f(x)=a_nx^n+a_{n-1}x^{n-1}+\cdots+a_1x+a_0,\quad a_n\neq 0,
\]
\[
g(x)=b_mx^m+b_{m-1}x^{m-1}+\cdots+b_1x+b_0,\quad b_m\neq 0,
\]
由于 $n\geq m$，可令
\[
f_1(x)=f(x)-a_nb_m^{-1}x^{n-m}g(x),
\]
则 $\deg f_1(x)<n$. 由归纳假设，有
\[
f_1(x)=g(x)q_1(x)+r(x),
\]
且 $\deg r(x)<\deg g(x)$，于是
\[
f(x)-a_nb_m^{-1}x^{n-m}g(x)=g(x)q_1(x)+r(x).
\]
因此
\[
f(x)=g(x)(a_nb_m^{-1}x^{n-m}+q_1(x))+r(x).
\]
令
\[
q(x)=a_nb_m^{-1}x^{n-m}+q_1(x),
\]
即得 (5.2.1) 式.

再证明唯一性. 设另有 $p(x),t(x)$，使
\[
f(x)=g(x)p(x)+t(x),
\]
且 $\deg t(x)<\deg g(x)$，则
\begin{equation}
g(x)(q(x)-p(x))=t(x)-r(x).
\tag{5.2.2}
\end{equation}
注意 (5.2.2) 式左边若 $q(x)-p(x)\neq 0$，便有
\[
\deg g(x)(q(x)-p(x))\geq \deg g(x)>\deg(t(x)-r(x)),
\]
引出矛盾. 因此只可能 $p(x)=q(x),\ t(x)=r(x).\quad \square$

\noindent\textbf{推论 5.2.1}\quad 设 $f(x),g(x)\in\mathbb{K}[x],\ g(x)\neq 0$，则 $g(x)\mid f(x)$ 的充分必要条件是 $g(x)$ 除 $f(x)$ 后的余式为零.

为了求出 (5.2.1) 式中的 $q(x)$ 及 $r(x)$，我们可用除法算式，参见下面的例子.

\noindent\textbf{例 5.2.1}\quad 设 $f(x)=3x^4-4x^3+5x-1,\ g(x)=x^2-x+1$，试求 (5.2.1) 式中的 $q(x)$ 及 $r(x)$.

\noindent\textbf{解}\quad 用长除法可得
\[
q(x)=3x^2-x-4,\qquad r(x)=2x+3.
\]

\subsection*{习题 5.2}

1. 设 $f(x)=2x^5+x^4-x+1,\ g(x)=x^3-x+2$，试求 $q(x)$ 及 $r(x)$，使
\[
f(x)=g(x)q(x)+r(x),
\]
且 $\deg r(x)<3$.

2. 设 $f(x)=2x^4-3x^3+4x^2+ax+b,\ g(x)=x^2-3x+1$. 若 $f(x)$ 除以 $g(x)$ 后余式为 $25x-5$，试求 $a,b$ 之值.

3. 设 $g(x)=ax+b\in\mathbb{K}[x]$ 且 $a\neq 0$，又 $f(x)\in\mathbb{K}[x]$，证明：$g(x)\mid f(x)^2$ 的充分必要条件是 $g(x)\mid f(x)$.

4. 若 $f(x)=x^4+px^2+q,\ g(x)=x^2+mx+1$，问：$p,q,m$ 适合什么条件才有 $g(x)\mid f(x)$？

5. 设 $g(x)=ax^2+bx+c(abc\neq 0),\ f(x)=x^3+px^2+qx+r$，满足 $g(x)\mid f(x)$，求证：
\[
\frac{ap-b}{a}=\frac{aq-c}{b}=\frac{ar}{c}.
\]

6. 证明：$x^d-a^d\mid x^n-a^n$ 的充分必要条件是 $d\mid n$，其中 $a\neq 0$.

7. 设 $f(x)=x^{3m}+x^{3n+1}+x^{3p+2}$，其中 $m,n,p$ 为自然数，又 $g(x)=x^2+x+1$，求证：$g(x)\mid f(x)$.

\section*{§5.3\quad 最大公因式}

设 $f(x),g(x)$ 是数域 $\mathbb{K}$ 上的多项式. 若 $h(x)\in\mathbb{K}[x]$ 适合：$h(x)\mid f(x),\ h(x)\mid g(x)$，则称 $h(x)$ 是 $f(x)$ 与 $g(x)$ 的公因式 (或公因子)；若 $l(x)\in\mathbb{K}[x]$ 适合：$f(x)\mid l(x),\ g(x)\mid l(x)$，则称 $l(x)$ 是 $f(x)$ 与 $g(x)$ 的公倍式.

\noindent\textbf{定义 5.3.1}\quad 设 $f(x),g(x)\in\mathbb{K}[x]$，若 $d(x)$ 是 $f(x)$ 与 $g(x)$ 的公因式，且对 $f(x)$ 与 $g(x)$ 的任一公因式 $h(x)$ 均有 $h(x)\mid d(x)$，则称 $d(x)$ 为 $f(x)$ 与 $g(x)$ 的最大公因式 (或称 $d(x)$ 为 $f(x),g(x)$ 的 g.c.d.)，记为 $d(x)=(f(x),g(x))$.

同理，若 $m(x)$ 是 $f(x)$ 与 $g(x)$ 的公倍式，且对 $f(x)$ 与 $g(x)$ 的任一公倍式 $l(x)$ 均有 $m(x)\mid l(x)$，则称 $m(x)$ 为 $f(x)$ 与 $g(x)$ 的最小公倍式 (或称 $m(x)$ 为 $f(x),g(x)$ 的 l.c.m.)，记为 $m(x)=[f(x),g(x)]$.

如何求两个多项式 $f(x),g(x)$ 的最大公因式 $d(x)$？不妨假设 $\deg f(x)\geq \deg g(x)$，由带余除法，有
\[
f(x)=g(x)q(x)+r(x),
\]
其中 $\deg r(x)<\deg g(x)$. 若 $r(x)\neq 0$，因为 $d(x)\mid f(x),d(x)\mid g(x)$，故 $d(x)\mid r(x)$，这表明 $d(x)$ 是 $g(x)$ 和 $r(x)$ 的公因式. 注意到 $r(x)$ 的次数比 $f(x),g(x)$ 都小. 如果我们再将 $g(x)$ 除以 $r(x)$，得到的余式次数更小 (如果不为零的话)，而 $d(x)$ 是 $r(x)$ 和这个余式的公因式. 不断这样做下去，肯定会得到余式为零的除式，其中的一个因式便是最大公因式. 这个方法称为 Euclid (欧几里得) 辗转相除法.

\noindent\textbf{定理 5.3.1}\quad 设 $f(x),g(x)\in\mathbb{K}[x]$，则 $f(x)$ 与 $g(x)$ 的最大公因式 $d(x)$ 必存在，且有 $u(x),v(x)\in\mathbb{K}[x]$，使
\begin{equation}
f(x)u(x)+g(x)v(x)=d(x).
\tag{5.3.1}
\end{equation}

\noindent\textbf{证明}\quad 若 $f(x)=0$，则显然 $(f(x),g(x))=g(x)$；若 $g(x)=0$，则 $(f(x),g(x))=f(x)$. 故不妨设 $f(x)\neq 0,\ g(x)\neq 0$. 由带余除法，我们有下列等式：
\[
\begin{aligned}
f(x)&=g(x)q_1(x)+r_1(x),\\
g(x)&=r_1(x)q_2(x)+r_2(x),\\
r_1(x)&=r_2(x)q_3(x)+r_3(x),\\
&\cdots\cdots\cdots\\
r_{s-2}(x)&=r_{s-1}(x)q_s(x)+r_s(x),\\
&\cdots\cdots\cdots
\end{aligned}
\]
余式的次数是严格递减的，因此经过有限步后，必有一个等式其余式为零. 不妨设 $r_{s+1}(x)=0$，于是
\[
r_{s-1}(x)=r_s(x)q_{s+1}(x).
\]
现在要证明 $r_s(x)$ 即为 $f(x)$ 与 $g(x)$ 的最大公因式. 由上式知 $r_s(x)\mid r_{s-1}(x)$，但
\begin{equation}
r_{s-2}(x)=r_{s-1}(x)q_s(x)+r_s(x),
\tag{5.3.2}
\end{equation}
因此 $r_s(x)\mid r_{s-2}(x)$. 这样可一直推下去，得到 $r_s(x)\mid g(x),\ r_s(x)\mid f(x)$. 这表明 $r_s(x)$ 是 $f(x)$ 与 $g(x)$ 的公因式. 又设 $h(x)$ 是 $f(x)$ 与 $g(x)$ 的公因式，则 $h(x)\mid r_1(x)$，于是 $h(x)\mid r_2(x)$，不断往下推，容易看出有 $h(x)\mid r_s(x)$. 因此 $r_s(x)$ 是最大公因式.

再证明 (5.3.1) 式. 从 (5.3.2) 式得
\begin{equation}
r_s(x)=r_{s-2}(x)-r_{s-1}(x)q_s(x),
\tag{5.3.3}
\end{equation}
但我们有
\begin{equation}
r_{s-3}(x)=r_{s-2}(x)q_{s-1}(x)+r_{s-1}(x),
\tag{5.3.4}
\end{equation}
从 (5.3.4) 式中解出 $r_{s-1}(x)$ 代入 (5.3.3) 式，得
\[
r_s(x)=r_{s-2}(x)(1+q_{s-1}(x)q_s(x))-r_{s-3}(x)q_s(x).
\]
用类似的方法逐步将 $r_i(x)$ 用 $r_{i-1}(x),r_{i-2}(x)$ 代入，最后得到
\[
r_s(x)=f(x)u(x)+g(x)v(x).
\]
显然 $u(x),v(x)\in\mathbb{K}[x].\quad \square$

若 $d_1(x)$ 与 $d_2(x)$ 都是 $f(x),g(x)$ 的最大公因式，则 $d_1(x)\mid d_2(x),\ d_2(x)\mid d_1(x)$. 因此必存在 $c\in\mathbb{K}$，使 $d_2(x)=cd_1(x)$，即 $f(x)$ 与 $g(x)$ 的两个最大公因式最多相差一个非零常数. 若规定 $f(x),g(x)$ 的最大公因式首项系数为 1 (首项系数等于 1 的多项式称为首一多项式)，则 $f(x)$ 与 $g(x)$ 的最大公因式就唯一确定了. 今后我们规定，凡最大公因式均指首一多项式.

对 $m$ 个多项式，也可以定义最大公因式的概念. 设 $f_i(x)(i=1,2,\cdots,m)$ 是 $\mathbb{K}[x]$ 中元素，若 $d(x)\mid f_i(x)(i=1,2,\cdots,m)$，则称 $d(x)$ 是 $\{f_i(x)\}$ 的公因式. 如果对 $\{f_i(x)\}$ 的任一公因式 $h(x),h(x)\mid d(x)$，则称 $d(x)$ 为 $\{f_i(x)\}$ 的最大公因式，记为 $d(x)=(f_1(x),f_2(x),\cdots,f_m(x))$.

\noindent\textbf{引理 5.3.1}\quad 设 $f_1(x),f_2(x),\cdots,f_m(x)\in\mathbb{K}[x]$，则
\[
((f_1(x),f_2(x)),f_3(x),\cdots,f_m(x))=(f_1(x),f_2(x),\cdots,f_m(x)).
\]

\noindent\textbf{证明}\quad 设 $d_1(x)=(f_1(x),f_2(x)),\ d(x)=(f_1(x),f_2(x),\cdots,f_m(x))$，则 $d(x)\mid d_1(x),\ d(x)\mid f_i(x)(i>2)$. 另一方面，若 $u(x)\mid d_1(x),\ u(x)\mid f_i(x)(i>2)$，则 $u(x)\mid f_1(x),\ u(x)\mid f_2(x)$，从而 $u(x)$ 是所有 $f_i(x)$ 的公因式，于是 $u(x)\mid d(x)$，故
\[
d(x)=(d_1(x),f_3(x),\cdots,f_m(x))=((f_1(x),f_2(x)),f_3(x),\cdots,f_m(x)).\quad \square
\]

最大公因式的定义显然与多项式的排列顺序无关，因此引理 5.3.1 告诉我们，求 $m$ 个多项式的最大公因式时可以先求其中任意两个的最大公因式，从而把问题化为 $m-1$ 个多项式的情形. 不断这样做下去，最后便可计算出 $m$ 个多项式的最大公因式.

\noindent\textbf{例 5.3.1}\quad $f(x)=x^4-x^3-x^2+2x-1,\ g(x)=x^3-2x+1$，求 $(f(x),g(x))$ 以及 $u(x),v(x)$，使 $f(x)u(x)+g(x)v(x)=(f(x),g(x))$.

\noindent\textbf{解}\quad 对 $f(x),g(x)$ 进行辗转相除，可得
\[
q_1(x)=x-1,\quad r_1(x)=x^2-x,\quad q_2(x)=x+1,\quad r_2(x)=-x+1,\quad q_3(x)=-x.
\]
因此
\[
(f(x),g(x))=x-1=-r_2(x).
\]
又
\[
\begin{aligned}
r_2(x)&=g(x)-r_1(x)q_2(x)\\
&=g(x)-(f(x)-g(x)q_1(x))q_2(x)\\
&=f(x)(-q_2(x))+g(x)(1+q_1(x)q_2(x)),
\end{aligned}
\]
所以
\[
u(x)=x+1,\qquad v(x)=-x^2.
\]

利用最大公因式，我们可以定义互素的概念.

\noindent\textbf{定义 5.3.2}\quad 设 $f(x),g(x)\in\mathbb{K}[x]$，若 $(f(x),g(x))=1$，则称 $f(x)$ 与 $g(x)$ 互素.

\noindent\textbf{定理 5.3.2}\quad 设 $f(x),g(x)\in\mathbb{K}[x]$，则 $f(x)$ 与 $g(x)$ 互素的充分必要条件是存在 $u(x),v(x)\in\mathbb{K}[x]$，使
\begin{equation}
f(x)u(x)+g(x)v(x)=1.
\tag{5.3.5}
\end{equation}

\noindent\textbf{证明}\quad 若 $(f(x),g(x))=1$，则由定理 5.3.1 可知，存在 $u(x),v(x)$，使 (5.3.5) 式成立. 反之，若 (5.3.5) 式成立，假设 $(f(x),g(x))=d(x)$，由 (5.3.5) 式可知 $d(x)\mid 1$，于是 $d(x)=1.\quad \square$

\noindent\textbf{推论 5.3.1}\quad 若 $f_1(x)\mid g(x),\ f_2(x)\mid g(x)$，且 $(f_1(x),f_2(x))=1$，则
\[
f_1(x)f_2(x)\mid g(x).
\]

\noindent\textbf{证明}\quad 设 $u(x),v(x)\in\mathbb{K}[x]$，使
\[
f_1(x)u(x)+f_2(x)v(x)=1.
\]
设 $g(x)=f_1(x)s(x)=f_2(x)t(x)$，则
\[
\begin{aligned}
g(x)&=g(x)(f_1(x)u(x)+f_2(x)v(x))\\
&=f_2(x)t(x)f_1(x)u(x)+f_1(x)s(x)f_2(x)v(x)\\
&=f_1(x)f_2(x)(t(x)u(x)+s(x)v(x)),
\end{aligned}
\]
即 $f_1(x)f_2(x)\mid g(x).\quad \square$

\noindent\textbf{推论 5.3.2}\quad 若 $(f(x),g(x))=1$，且 $f(x)\mid g(x)h(x)$，则
\[
f(x)\mid h(x).
\]

\noindent\textbf{证明}\quad 设 $u(x),v(x)\in\mathbb{K}[x]$，使
\[
f(x)u(x)+g(x)v(x)=1,
\]
则
\[
f(x)u(x)h(x)+g(x)v(x)h(x)=h(x).
\]
因上式左边可被 $f(x)$ 整除，故 $f(x)\mid h(x).\quad \square$

\noindent\textbf{推论 5.3.3}\quad 设 $(f(x),g(x))=d(x),\ f(x)=f_1(x)d(x),\ g(x)=g_1(x)d(x)$，则
\[
(f_1(x),g_1(x))=1.
\]

\noindent\textbf{证明}\quad 设 $u(x),v(x)\in\mathbb{K}[x]$，使
\[
f(x)u(x)+g(x)v(x)=d(x),
\]
即
\[
f_1(x)d(x)u(x)+g_1(x)d(x)v(x)=d(x),
\]
两边消去 $d(x)$ 即得
\[
f_1(x)u(x)+g_1(x)v(x)=1,
\]
因此 $f_1(x),g_1(x)$ 互素. $\square$

\noindent\textbf{推论 5.3.4}\quad 设 $(f(x),g(x))=d(x)$，则
\[
(t(x)f(x),t(x)g(x))=t(x)d(x).
\]

\noindent\textbf{证明}\quad 设 $u(x),v(x)\in\mathbb{K}[x]$，使
\[
f(x)u(x)+g(x)v(x)=d(x),
\]
则
\[
t(x)f(x)u(x)+t(x)g(x)v(x)=t(x)d(x).
\]
因此，若 $h(x)\mid t(x)f(x),\ h(x)\mid t(x)g(x)$，则必有 $h(x)\mid t(x)d(x)$. 又 $t(x)d(x)$ 是 $t(x)f(x),t(x)g(x)$ 的公因式，因此 $t(x)d(x)$ 是 $t(x)f(x)$ 与 $t(x)g(x)$ 的最大公因式. $\square$

\noindent\textbf{推论 5.3.5}\quad 若 $(f_1(x),g(x))=1,\ (f_2(x),g(x))=1$，则
\[
(f_1(x)f_2(x),g(x))=1.
\]

\noindent\textbf{证明}\quad 设
\[
f_1(x)u_1(x)+g(x)v_1(x)=1,
\]
\[
f_2(x)u_2(x)+g(x)v_2(x)=1,
\]
将上两式两边分别相乘得
\[
f_1(x)f_2(x)u_1(x)u_2(x)+g(x)(v_1(x)f_2(x)u_2(x)+g(x)v_1(x)v_2(x)+v_2(x)f_1(x)u_1(x))=1.
\]
这就是说 $f_1(x)f_2(x)$ 和 $g(x)$ 互素. $\square$

下面的推论证明了最小公倍式的存在性. 与最大公因式类似的讨论可知，$f(x)$ 与 $g(x)$ 的两个最小公倍式最多相差一个非零常数. 今后我们规定，凡最小公倍式均指首一多项式，这就保证了最小公倍式的唯一性.

\noindent\textbf{推论 5.3.6}\quad 设 $f(x),g(x)$ 是非零多项式，则
\[
f(x)g(x)\sim (f(x),g(x))[f(x),g(x)].
\]

\noindent\textbf{证明}\quad 设 $d(x)=(f(x),g(x))$ 且 $f(x)=f_0(x)d(x),\ g(x)=g_0(x)d(x)$，则 $f_0(x),g_0(x)$ 互素. 假设 $l(x)$ 是 $f(x),g(x)$ 的公倍式且
\[
l(x)=f(x)u(x)=g(x)v(x),
\]
则 $f_0(x)d(x)u(x)=g_0(x)d(x)v(x)$，消去 $d(x)$ 得
\[
f_0(x)u(x)=g_0(x)v(x).
\]
因为 $f_0(x),g_0(x)$ 互素，所以 $f_0(x)\mid v(x),\ g_0(x)\mid u(x)$. 设 $u(x)=g_0(x)p(x)$，则
\[
l(x)=f_0(x)d(x)g_0(x)p(x),
\]
即 $f_0(x)d(x)g_0(x)\mid l(x)$. 显然 $f_0(x)d(x)g_0(x)$ 是 $f(x),g(x)$ 的公倍式，因此它与 $f(x),g(x)$ 的最小公倍式相差一个非零常数，即
\[
\frac{f(x)g(x)}{d(x)}=f_0(x)d(x)g_0(x)\sim [f(x),g(x)].\quad \square
\]

下面的定理是多项式形式的“中国剩余定理”. “中国剩余定理”又称为“孙子定理”，是一个在数论、代数学等领域有广泛应用的定理.

\noindent\textbf{定理 5.3.3 (中国剩余定理)}\quad 设 $g_1(x),\cdots,g_n(x)$ 是两两互素的多项式，$r_1(x),\cdots,r_n(x)$ 是 $n$ 个多项式，则存在多项式 $f(x),q_1(x),\cdots,q_n(x)$，使
\[
f(x)=g_i(x)q_i(x)+r_i(x),\quad i=1,\cdots,n.
\]

\noindent\textbf{证明}\quad 先证明存在多项式 $f_i(x)$，使对任意的 $i$，有
\[
f_i(x)=g_i(x)p_i(x)+1,\quad g_j(x)\mid f_i(x)\ (j\neq i).
\]
一旦得证，只需令 $f(x)=r_1(x)f_1(x)+\cdots+r_n(x)f_n(x)$ 即可. 现构造 $f_1(x)$ 如下. 因为 $g_1(x)$ 和 $g_j(x)(j\neq 1)$ 互素，故存在 $u_j(x),v_j(x)$，使 $g_1(x)u_j(x)+g_j(x)v_j(x)=1$. 令
\[
f_1(x)=g_2(x)v_2(x)\cdots g_n(x)v_n(x)=(1-g_1(x)u_2(x))\cdots(1-g_1(x)u_n(x)),
\]
显然 $f_1(x)$ 符合要求. 同理可构造 $f_i(x).\quad \square$

\subsection*{习题 5.3}

1. 用辗转相除法求下列多项式的最大公因式：

(1) $f(x)=x^4-x^3-4x^2-x+1,\ g(x)=x^2-x-2$；

(2) $f(x)=x^4+3x^3-2x-2,\ g(x)=x^2+2x-3$；

(3) $f(x)=x^5-2x^4-4x^3+2x^2+4x+6,\ g(x)=x^2+x+1$.

2. 求 $u(x),v(x)$，使 $f(x)u(x)+g(x)v(x)=(f(x),g(x))$：

(1) $f(x)=x^4+2x^3-x^2-4x-2,\ g(x)=x^4+x^3-x^2-2x-2$；

(2) $f(x)=x^4-x^3-4x^2+4x+1,\ g(x)=x^2-x-1$.

3. 设 $d(x)=f(x)u(x)+g(x)v(x)$，举例说明 $d(x)$ 不必是 $f(x)$ 与 $g(x)$ 的最大公因式. 若进一步有 $d(x)\mid f(x),d(x)\mid g(x)$，求证：$d(x)$ 必是 $f(x)$ 与 $g(x)$ 的最大公因式.

4. 设 $f(x)$ 与 $g(x)$ 互素，求证：$f(x^m)$ 与 $g(x^m)$ 也互素，其中 $m$ 为任一正整数.

5. 设 $(f(x),g(x))=1$，求证：$(f(x)g(x),f(x)+g(x))=1$.

6. 设 $f_1(x),f_2(x),\cdots,f_m(x),g_1(x),g_2(x),\cdots,g_n(x)$ 为多项式，且
\[
(f_i(x),g_j(x))=1,\quad i=1,\cdots,m;\ j=1,\cdots,n,
\]
求证：$(f_1(x)f_2(x)\cdots f_m(x),g_1(x)g_2(x)\cdots g_n(x))=1$.

7. 设 $f_1(x),f_2(x),\cdots,f_n(x)$ 的最大公因式为 $d(x)$，求证：必存在 $g_1(x),g_2(x),\cdots,g_n(x)$，使
\[
f_1(x)g_1(x)+f_2(x)g_2(x)+\cdots+f_n(x)g_n(x)=d(x).
\]
特别，$f_1(x),f_2(x),\cdots,f_n(x)$ 互素 $(d(x)=1)$ 的充分必要条件是存在 $g_1(x),g_2(x),\cdots,g_n(x)$，使
\[
f_1(x)g_1(x)+f_2(x)g_2(x)+\cdots+f_n(x)g_n(x)=1.
\]

8. 设 $f_1(x),f_2(x),\cdots,f_m(x)\in\mathbb{K}[x]$，若 $f_i(x)\mid m(x)(i=1,2,\cdots,m)$，则称 $m(x)$ 是 $\{f_i(x)\}$ 的公倍式. 进一步，若对 $\{f_i(x)\}$ 的任一公倍式 $l(x),m(x)\mid l(x)$，则称 $m(x)$ 为 $\{f_i(x)\}$ 的最小公倍式，记为 $m(x)=[f_1(x),f_2(x),\cdots,f_m(x)]$. 求证：
\[
[[f_1(x),f_2(x)],f_3(x),\cdots,f_m(x)]=[f_1(x),f_2(x),\cdots,f_m(x)].
\]

\section*{§5.4\quad 因式分解}

读者在中学里已经学习过因式分解. 但是由于当时没有数域的概念，因此对什么叫“不可再分”缺乏严格的定义. 现在我们可以来严格地建立多项式的因式分解理论.

\noindent\textbf{定义 5.4.1}\quad 设 $f(x)$ 是数域 $\mathbb{K}$ 上的非常数多项式，若 $f(x)$ 可以分解为两个次数小于 $f(x)$ 次数的 $\mathbb{K}$ 上多项式之积，则称 $f(x)$ 是 $\mathbb{K}$ 上的可约多项式. 否则，称 $f(x)$ 为 $\mathbb{K}$ 上的不可约多项式.

从这个定义看出，多项式的可约或不可约与数域密切相关：比如 $x^2-2$ 在有理数域上是不可约多项式，但在实数域上是可约多项式. 然而无论在哪个域上，一次多项式总是不可约的. 另外我们还要注意，因为 $\mathbb{K}$ 是一个域，对 $\mathbb{K}$ 中任一非零数 $c$，都有 $f(x)=c(c^{-1}f(x))$，但是这种分解没有多大意义. 我们谈的因式分解，都是指将一个多项式分解为两个次数较小的多项式之积.

\noindent\textbf{引理 5.4.1}\quad 设 $f(x)$ 是数域 $\mathbb{K}$ 上的不可约多项式，则对 $\mathbb{K}$ 上任一多项式 $g(x)$，或者 $f(x)\mid g(x)$，或者 $(f(x),g(x))=1$.

\noindent\textbf{证明}\quad 设 $d(x)=(f(x),g(x))$. 因为 $f(x)$ 不可约，故 $f(x)$ 的因式只能是非零常数多项式或 $cf(x)(c\neq 0)$，从而或者 $d(x)=1$ 或者 $d(x)=cf(x)$ (首一多项式)，故得结论. $\square$

不可约多项式还有“素性”，如以下定理所述.

\noindent\textbf{定理 5.4.1}\quad 设 $p(x)$ 是 $\mathbb{K}$ 上的不可约多项式，$f(x),g(x)$ 是 $\mathbb{K}$ 上的多项式且 $p(x)\mid f(x)g(x)$，则或者 $p(x)\mid f(x)$，或者 $p(x)\mid g(x)$.

\noindent\textbf{证明}\quad 若 $p(x)$ 不能整除 $f(x)$，则由引理 5.4.1 知 $p(x)$ 与 $f(x)$ 互素，再由推论 5.3.2 即得 $p(x)\mid g(x).\quad \square$

\noindent\textbf{推论 5.4.1}\quad 设 $p(x)$ 为不可约多项式且
\[
p(x)\mid f_1(x)f_2(x)\cdots f_m(x),
\]
则 $p(x)$ 必可整除其中某个 $f_i(x)$.

对 $\mathbb{K}[x]$ 中任意一个非常数多项式，是否一定可以分解为不可约因子的积？这种分解是否唯一？下面的因式分解定理回答了这个问题.

\noindent\textbf{定理 5.4.2}\quad 设 $f(x)$ 是数域 $\mathbb{K}$ 上的多项式且 $\deg f(x)\geq 1$，则

(1) $f(x)$ 可分解为有限个 $\mathbb{K}$ 上的不可约多项式之积；

(2) 若
\begin{equation}
f(x)=p_1(x)p_2(x)\cdots p_s(x)=q_1(x)q_2(x)\cdots q_t(x)
\tag{5.4.1}
\end{equation}
是 $f(x)$ 的两个不可约分解，即 $p_i(x),q_j(x)$ 都是 $\mathbb{K}$ 上的次数大于零的不可约多项式，则 $s=t$，且经过适当调换因式的次序以后，有
\[
q_i(x)\sim p_i(x),\quad i=1,2,\cdots,s.
\]

\noindent\textbf{证明}\quad (1) 对多项式 $f(x)$ 的次数用数学归纳法. 若 $\deg f(x)=1$，结论显然成立. 设次数小于 $n$ 的多项式都可以分解为 $\mathbb{K}$ 上的不可约多项式之积而 $\deg f(x)=n$. 若 $f(x)$ 不可约，结论自然成立. 若 $f(x)$ 可约，则
\[
f(x)=f_1(x)f_2(x),
\]
其中 $f_1(x),f_2(x)$ 的次数小于 $n$，由归纳假设它们可以分解为有限个 $\mathbb{K}$ 上的不可约多项式之积. 所有这些多项式之积就是 $f(x)$.

(2) 对 (5.4.1) 式中的 $s$ 用数学归纳法. 若 $s=1$，则 $f(x)=p_1(x)$，因此 $f(x)$ 是不可约多项式，于是 $t=1,\ q_1(x)=p_1(x)$. 现假设对不可约因子个数小于 $s$ 的多项式结论正确. 由 (5.4.1) 式，有
\[
p_1(x)\mid q_1(x)q_2(x)\cdots q_t(x),
\]
由推论 5.4.1 可知，必存在某个 $i$，不妨设 $i=1$，使
\[
p_1(x)\mid q_1(x).
\]
但是 $p_1(x),q_1(x)$ 都是不可约多项式，因此存在 $0\neq c_1\in\mathbb{K}$，使
\[
q_1(x)=c_1p_1(x),
\]
此即 $p_1(x)\sim q_1(x)$. 将上式代入 (5.4.1) 式并消去 $p_1(x)$，得到
\[
p_2(x)\cdots p_s(x)=c_1q_2(x)\cdots q_t(x).
\]
这时左边为 $s-1$ 个不可约多项式之积，由归纳假设，$s-1=t-1$，即 $s=t$. 另一方面，存在 $0\neq c_i\in\mathbb{K}$，使 $q_i(x)=c_ip_i(x).\quad \square$

定理 5.4.2 表明，任一多项式可唯一地分解为若干个不可约多项式之积. 这里唯一是在相伴意义下的唯一，即相应的多项式可以差一个常数因子. 如果把分解式中相同或仅差一个常数的因式合并在一起，就得到了一个“标准分解”式：
\begin{equation}
f(x)=cp_1(x)^{e_1}p_2(x)^{e_2}\cdots p_m(x)^{e_m},
\tag{5.4.2}
\end{equation}
其中 $c\neq 0$，$p_i(x)$ 是互异的首一不可约多项式，$e_i\geq 1(i=1,2,\cdots,m)$.

若 $e_i>1(e_i=1)$，我们称 (5.4.2) 式中的因式 $p_i(x)$ 为 $f(x)$ 的 $e_i$ 重因式 (单因式). 显然这时 $p_i(x)^{e_i}\mid f(x)$，但 $p_i(x)^{e_i+1}$ 不能整除 $f(x)$.

\noindent\textbf{例 5.4.1}\quad 设 $f(x),g(x)$ 是 $\mathbb{K}$ 上的两个多项式，在它们的标准分解式中适当添加零次项，故不妨设它们有如下的分解式：
\[
f(x)=c_1p_1(x)^{e_1}p_2(x)^{e_2}\cdots p_n(x)^{e_n},
\]
\[
g(x)=c_2p_1(x)^{f_1}p_2(x)^{f_2}\cdots p_n(x)^{f_n},
\]
其中 $e_i\geq 0,\ f_i\geq 0(i=1,2,\cdots,n)$，则 $f(x),g(x)$ 的最大公因式
\[
(f(x),g(x))=p_1(x)^{k_1}p_2(x)^{k_2}\cdots p_n(x)^{k_n},
\]
其中 $k_i=\min\{e_i,f_i\}(i=1,2,\cdots,n)$.

类似地，$f(x),g(x)$ 的最小公倍式
\[
[f(x),g(x)]=p_1(x)^{h_1}p_2(x)^{h_2}\cdots p_n(x)^{h_n},
\]
其中 $h_i=\max\{e_i,f_i\}(i=1,2,\cdots,n)$.

一个多项式何时有重因式？我们现在来讨论这个问题. 为此，首先我们回忆一下数学分析中导数的概念.

若 $f(x)=a_nx^n+a_{n-1}x^{n-1}+\cdots+a_1x+a_0$，则 $f(x)$ 的导数或微商为下列多项式：
\begin{equation}
f'(x)=na_nx^{n-1}+(n-1)a_{n-1}x^{n-2}+\cdots+a_1.
\tag{5.4.3}
\end{equation}
事实上，我们也可以直接定义 (5.4.3) 式为多项式 $f(x)$ 的形式导数而不借用数学分析中的定义. 按此定义不难验证下列公式：

(1) $(f(x)+g(x))'=f'(x)+g'(x)$；

(2) $(cf(x))'=cf'(x)$；

(3) $(f(x)g(x))'=f'(x)g(x)+f(x)g'(x)$；

(4) $(f(x)^m)'=mf(x)^{m-1}f'(x)$.

\noindent\textbf{定理 5.4.3}\quad 数域 $\mathbb{K}$ 上的多项式 $f(x)$ 没有重因式的充分必要条件是 $f(x)$ 与 $f'(x)$ 互素.

\noindent\textbf{证明}\quad 设多项式 $p(x)$ 是 $f(x)$ 的 $m(m>1)$ 重因式，则 $f(x)=p(x)^mg(x)$，故
\[
f'(x)=mp(x)^{m-1}p'(x)g(x)+p(x)^mg'(x).
\]
于是 $p(x)^{m-1}\mid f'(x)$，这表明 $f(x)$ 与 $f'(x)$ 有公因式 $p(x)^{m-1}$.

反之，若不可约多项式 $p(x)$ 是 $f(x)$ 的单因式，可设 $f(x)=p(x)g(x),p(x)$ 不能整除 $g(x)$. 于是
\[
f'(x)=p'(x)g(x)+p(x)g'(x).
\]
若 $p(x)$ 是 $f'(x)$ 的因式，则 $p(x)\mid p'(x)g(x)$. 但 $p(x)$ 不能整除 $g(x)$ 且 $p(x)$ 不可约，故 $p(x)\mid p'(x)$. 而 $p'(x)\neq 0$ 且 $\deg p'(x)<\deg p(x)$，这是不可能的. 若 $f(x)$ 无重因式，则在 $f(x)$ 的标准分解式 (5.4.2) 中，$e_i=1$ 对一切 $i=1,2,\cdots,m$ 成立，于是 $p_i(x)$ 都不能整除 $f'(x)$. 由于 $p_i(x)$ 为不可约多项式，故 $(p_i(x),f'(x))=1$，由推论 5.3.5 可知
\[
(p_1(x)p_2(x)\cdots p_m(x),f'(x))=1,
\]
即 $(f(x),f'(x))=1.\quad \square$

定理 5.4.3 提供了判断多项式 $f(x)$ 是否有重因式的方法：用 Euclid 辗转相除法求出 $f(x)$ 及其导数 $f'(x)$ 的最大公因式. 若 $(f(x),f'(x))=1$，则 $f(x)$ 无重因式，否则必有重因式. 这个方法的好处在于它不必求出 $f(x)$ 的标准分解式，也不必求根即可作出判断 (事实上，求标准分解式或求根往往非常困难).

当 $f(x)$ 有重因式时，我们可借助 $(f(x),f'(x))$ 将 $f(x)$ 的重因式消去，即有下列命题.

\noindent\textbf{命题 5.4.1}\quad 设 $d(x)=(f(x),f'(x))$，则 $f(x)/d(x)$ 是一个没有重因式的多项式，且这个多项式的不可约因式与 $f(x)$ 的不可约因式相同 (不计重数).

\noindent\textbf{证明}\quad 设 $f(x)$ 有如 (5.4.2) 式的标准分解式，则
\begin{equation}
\begin{aligned}
f'(x)=&\ ce_1p_1(x)^{e_1-1}p_2(x)^{e_2}\cdots p_s(x)^{e_s}p_1'(x)\\
&+ce_2p_1(x)^{e_1}p_2(x)^{e_2-1}\cdots p_s(x)^{e_s}p_2'(x)\\
&+\cdots\\
&+ce_sp_1(x)^{e_1}p_2(x)^{e_2}\cdots p_s(x)^{e_s-1}p_s'(x).
\end{aligned}
\tag{5.4.4}
\end{equation}
因此 $p_1(x)^{e_1-1}p_2(x)^{e_2-1}\cdots p_s(x)^{e_s-1}$ 是 $f(x)$ 与 $f'(x)$ 的公因式. 注意到 $f(x)$ 的因式一定具有 $p_1(x)^{k_1}p_2(x)^{k_2}\cdots p_s(x)^{k_s}$ 的形状. 不妨设 $h(x)$ 是 $f(x),f'(x)$ 的公因式. 注意到 $p_1(x)^{e_1}$ 可以整除 (5.4.4) 式中右边除第一项外的所有项，但不能整除第一项，因此 $p_1(x)^{e_1}$ 不能整除 $f'(x)$. 同理，$p_i(x)^{e_i}$ 不能整除 $f'(x)$. 由此我们不难看出
\[
h(x)\mid p_1(x)^{e_1-1}p_2(x)^{e_2-1}\cdots p_s(x)^{e_s-1},
\]
即 $p_1(x)^{e_1-1}p_2(x)^{e_2-1}\cdots p_s(x)^{e_s-1}=d(x)$. 显然 $f(x)/d(x)$ 没有重因式且与 $f(x)$ 含有相同的不可约因式. $\square$

\subsection*{习题 5.4}

1. 判断下列多项式有无重因式：

(1) $f(x)=x^5-10x^3-20x^2-15x-4$；

(2) $f(x)=4x^3-4x^2-7x-2$；

(3) $f(x)=x^3+x+1$.

2. 设 $\mathbb{K}_1\subseteq\mathbb{K}_2$ 是两个数域，若 $p(x)\in\mathbb{K}_1[x]$ 是 $\mathbb{K}_2$ 上的不可约多项式，求证：$p(x)$ 在 $\mathbb{K}_1$ 上也不可约.

3. 设 $p(x)$ 是数域 $\mathbb{K}$ 上的非常数多项式，求证：$p(x)$ 为 $\mathbb{K}$ 上不可约多项式的充分必要条件是对 $\mathbb{K}$ 上任意适合 $p(x)\mid f(x)g(x)$ 的多项式 $f(x)$ 与 $g(x)$，或者 $p(x)\mid f(x)$，或者 $p(x)\mid g(x)$.

4. 设 $u$ 是复数域中某个数，若 $u$ 适合某个非零有理系数多项式 (或整系数多项式) $f(x)=a_nx^n+a_{n-1}x^{n-1}+\cdots+a_0$，即 $f(u)=0$，则称 $u$ 是一个代数数. 证明：

(1) 对任一代数数 $u$，存在唯一一个 $u$ 适合的首一有理系数多项式 $g(x)$，使 $g(x)$ 是 $u$ 适合的所有非零有理系数多项式中的次数最小者. 这样的 $g(x)$ 称为 $u$ 的极小多项式或最小多项式.

(2) 设 $g(x)$ 是一个 $u$ 适合的首一有理系数多项式，则 $g(x)$ 是 $u$ 的极小多项式的充分必要条件是 $g(x)$ 是有理数域上的不可约多项式.

5. 证明：$g(x)^2\mid f(x)^2$ 的充分必要条件是 $g(x)\mid f(x)$.

6. 设 $f(x),g(x)$ 都是数域 $\mathbb{K}$ 上的多项式，$n$ 是一个正整数，证明：
\[
(f(x)^n,g(x)^n)=(f(x),g(x))^n,\quad [f(x)^n,g(x)^n]=[f(x),g(x)]^n.
\]

\section*{§5.5\quad 多项式函数}

我们在定义数域 $\mathbb{K}$ 上的多项式 $f(x)$ 时，未定元 $x$ 被看成是一个形式元，多项式 $f(x)$ 是一个形式多项式. 若
\[
f(x)=a_nx^n+a_{n-1}x^{n-1}+\cdots+a_1x+a_0,
\]
对 $\mathbb{K}$ 中任一元 $b$，定义
\[
f(b)=a_nb^n+a_{n-1}b^{n-1}+\cdots+a_1b+a_0,
\]
则称 $f(b)$ 为 $f(x)$ 在点 $b$ 的值. 这样多项式 $f(x)$ 又可看成是数域 $\mathbb{K}$ 上的函数，这个函数称为多项式函数. 一个很自然的问题是：多项式函数与多项式是否是一回事？也就是说，若两个多项式 $f(x)$ 与 $g(x)$ 在 $\mathbb{K}$ 上取值相同，那么是否必有 $f(x)=g(x)$，即它们对应的各次项的系数相同？我们将在下面给予肯定的回答.

\noindent\textbf{定义 5.5.1}\quad 设 $f(x)\in\mathbb{K}[x],\ b\in\mathbb{K}$，若 $f(b)=0$，则称 $b$ 是 $f(x)$ 的一个根或零点.

\noindent\textbf{定理 5.5.1 (余数定理)}\quad 设 $f(x)\in\mathbb{K}[x],\ b\in\mathbb{K}$，则存在 $g(x)\in\mathbb{K}[x]$，使
\[
f(x)=(x-b)g(x)+f(b).
\]
特别，$b$ 是 $f(x)$ 的根当且仅当 $(x-b)\mid f(x)$.

\noindent\textbf{证明}\quad 由带余除法知
\begin{equation}
f(x)=(x-b)g(x)+r(x),
\tag{5.5.1}
\end{equation}
其中 $\deg r(x)<1$，因此 $r(x)$ 为常数多项式. 在 (5.5.1) 式中用 $b$ 代替 $x$，即得 $r(x)=f(b).\quad \square$

由定理 5.5.1，我们可以仿照标准因式分解中重因式的概念来给出多项式重根的定义.

\noindent\textbf{定义 5.5.2}\quad 设 $f(x)\in\mathbb{K}[x],\ b\in\mathbb{K}$，若存在正整数 $k$，使 $(x-b)^k\mid f(x)$，但 $(x-b)^{k+1}$ 不能整除 $f(x)$，则称 $b$ 是 $f(x)$ 的一个 $k$ 重根. 若 $k=1$，则称 $b$ 为单根.

\noindent\textbf{引理 5.5.1}\quad 设 $f(x)$ 是数域 $\mathbb{K}$ 上的不可约多项式且 $\deg f(x)\geq 2$，则 $f(x)$ 在 $\mathbb{K}$ 中没有根.

\noindent\textbf{证明}\quad 用反证法，设 $b\in\mathbb{K}$ 是 $f(x)$ 的根，由定理 5.5.1 知 $(x-b)\mid f(x)$，即 $f(x)=(x-b)g(x)$ 可分解为两个低次多项式之积，这与 $f(x)$ 不可约矛盾. $\square$

如果把 $k$ 重根看成 $f(x)$ 有 $k$ 个根，则有下列结论.

\noindent\textbf{定理 5.5.2}\quad 若 $f(x)$ 是数域 $\mathbb{K}$ 上的 $n$ 次多项式，则 $f(x)$ 在 $\mathbb{K}$ 中最多只有 $n$ 个根.

\noindent\textbf{证明}\quad 将 $f(x)$ 作标准因式分解，则由引理 5.5.1 知 $f(x)$ 在 $\mathbb{K}$ 中根的个数等于该分解式中一次因式的个数，它不会超过 $n.\quad \square$

\noindent\textbf{推论 5.5.1}\quad 设 $f(x)$ 与 $g(x)$ 是 $\mathbb{K}$ 上的次数不超过 $n$ 的两个多项式，若存在 $\mathbb{K}$ 上 $n+1$ 个不同的数 $b_1,b_2,\cdots,b_{n+1}$，使
\[
f(b_i)=g(b_i),\quad i=1,2,\cdots,n+1,
\]
则 $f(x)=g(x)$.

\noindent\textbf{证明}\quad 作 $h(x)=f(x)-g(x)$，显然 $h(x)$ 次数不超过 $n$. 但它有 $n+1$ 个不同的根，因此只可能 $h(x)=0$，即 $f(x)=g(x).\quad \square$

推论 5.5.1 肯定地回答了我们一开始提出的问题，即若两个多项式 $f(x),g(x)$ 在数域 $\mathbb{K}$ 上的值相同，则 $f(x)=g(x)$. 但是读者必须注意，这是因为任一数域都有无限个元素之故. 对一般的域 (读者将在抽象代数教程中遇到) 来说，上述结论是不一定成立的，即存在两个不同的多项式，它们在某一有限域上的值处处相同.

\subsection*{习题 5.5}

1. 用余数定理证明：

(1) $x(x+1)(2x+1)\mid ((x+1)^{2n}-x^{2n}-2x-1)$；

(2) $(x^4+x^3+x^2+x+1)\mid (x^5+x^{11}+x^{17}+x^{23}+x^{29})$.

2. 证明：数域 $\mathbb{K}$ 上任意一个不可约多项式在复数域 $\mathbb{C}$ 中无重根.

3. 求证：$a$ 是多项式 $f(x)$ 的 $k$ 重根的充分必要条件是
\[
f(a)=f'(a)=\cdots=f^{(k-1)}(a)=0,\quad f^{(k)}(a)\neq 0.
\]

4. 设 $\deg f(x)=n\geq 1$，若 $f'(x)\mid f(x)$，证明：$f(x)$ 有 $n$ 重根.

5. 设 $f(x)$ 是数域 $\mathbb{K}$ 上的多项式，若对 $\mathbb{K}$ 中某个非零常数 $a$，有 $f(x+a)=f(x)$，求证：$f(x)$ 必是常数多项式.

6. 证明下列 Lagrange 插值定理：设 $a_1,a_2,\cdots,a_n$ 是 $\mathbb{K}$ 中 $n$ 个不同的数，$b_1,b_2,\cdots,b_n$ 是 $\mathbb{K}$ 中任意 $n$ 个数，则存在唯一的 $\mathbb{K}$ 上的多项式 $f(x),\deg f(x)\leq n-1$，且
\[
f(a_i)=b_i,\quad i=1,2,\cdots,n.
\]
试求出这个多项式.

7. 设 $f(x)$ 是一个 $n$ 次多项式，若当 $k=0,1,\cdots,n$ 时有 $f(k)=\dfrac{k}{k+1}$，求 $f(n+1)$.

8. 设 $(x^4+x^3+x^2+x+1)\mid (x^3f_1(x^5)+x^2f_2(x^5)+xf_3(x^5)+f_4(x^5))$，这里 $f_i(x)(i=1,2,3,4)$ 都是实系数多项式，求证：$f_i(1)=0(i=1,2,3,4)$.

\section*{§5.6\quad 复系数多项式}

现在我们来研究复数域上的多项式，首先我们要证明重要的“代数基本定理”. 这个定理的内容读者可能在中学里就已经知道. 定理的证明要用到二元连续函数的性质，对这方面不熟悉的读者可以跳过这一证明，承认其结论就可以了. 这个定理在复变函数课程里还将有一个非常简洁的证明. 在下面的证明里，我们不再区别多项式和多项式函数.

\noindent\textbf{定理 5.6.1 (代数基本定理)}\quad 次数大于零的复数域上的一元多项式至少有一个复数根.

\noindent\textbf{证明}\quad 设复数域上的 $n$ 次多项式为
\begin{equation}
f(z)=a_nz^n+a_{n-1}z^{n-1}+\cdots+a_1z+a_0.
\tag{5.6.1}
\end{equation}
首先证明，必存在一个复数 $z_0$，使对一切复数 $z$，有
\[
|f(z)|\geq |f(z_0)|.
\]
令 $z=x+iy$，其中 $x,y$ 是实变量. 展开 $f(x+iy)$ 并分开实部和虚部，则
\[
f(z)=u(x,y)+iv(x,y),
\]
其中 $u(x,y)$ 及 $v(x,y)$ 为实系数二元多项式函数. 又
\[
|f(z)|=\sqrt{u(x,y)^2+v(x,y)^2}
\]
是一个二元连续函数，但
\[
\begin{aligned}
|f(z)|&=|a_nz^n+a_{n-1}z^{n-1}+\cdots+a_0|\\
&\geq |a_nz^n|-|a_{n-1}z^{n-1}+\cdots+a_0|\\
&\geq |z|^n\left(|a_n|-\left(\frac{|a_{n-1}|}{|z|}+\frac{|a_{n-2}|}{|z|^2}+\cdots+\frac{|a_0|}{|z|^n}\right)\right),
\end{aligned}
\]
因此当 $|z|\to\infty$ 时，$|f(z)|\to\infty$. 于是必存在一个实常数 $R$，当 $|z|>R$ 时，$|f(z)|$ 充分大，因此 $|f(z)|$ 的最小值必含于圆 $|z|\leq R$ 中. 但这是平面上的一个闭区域，因此必存在 $z_0$，使 $|f(z_0)|$ 为最小.

接下去要证明 $f(z_0)=0$. 用反证法，即若 $f(z_0)\neq 0$，则必可找到 $z_1$，使 $|f(z_1)|<|f(z_0)|$，这样就与 $|f(z_0)|$ 是最小值相矛盾.

将 $z=z_0+h$ 代入 (5.6.1) 式便可得到一个关于 $h$ 的 $n$ 次多项式：
\begin{equation}
f(z_0+h)=b_nh^n+b_{n-1}h^{n-1}+\cdots+b_1h+b_0.
\tag{5.6.2}
\end{equation}
当 $h=0$ 时，$f(z_0)=b_0$，由假设 $f(z_0)\neq 0$，故
\[
\frac{f(z_0+h)}{f(z_0)}=\frac{b_n}{f(z_0)}h^n+\frac{b_{n-1}}{f(z_0)}h^{n-1}+\cdots+\frac{b_1}{f(z_0)}h+1.
\]
$b_1,b_2,\cdots,b_n$ 中有些可能为零，但绝不会全为零. 设 $b_k$ 是第一个不为零的复数，则
\begin{equation}
\frac{f(z_0+h)}{f(z_0)}=1+c_kh^k+c_{k+1}h^{k+1}+\cdots+c_nh^n,
\tag{5.6.3}
\end{equation}
其中 $c_j=\dfrac{b_j}{f(z_0)}$. 令 $d=\sqrt[k]{-\dfrac{1}{c_k}},\ h=ed$ 代入 (5.6.3) 式得
\[
\frac{f(z_0+h)}{f(z_0)}=1-e^k+e^{k+1}(c_{k+1}d^{k+1}+c_{k+2}d^{k+2}e+\cdots),
\]
取充分小的正实数 $e$ (至少小于 1)，使
\[
e(|c_{k+1}d^{k+1}|+|c_{k+2}d^{k+2}|+\cdots)<\frac12,
\]
于是
\[
\begin{aligned}
\left|\frac{f(z_0+h)}{f(z_0)}\right|
&\leq |1-e^k|+|e^{k+1}(c_{k+1}d^{k+1}+c_{k+2}d^{k+2}e+\cdots)|\\
&\leq 1-e^k+e^{k+1}(|c_{k+1}d^{k+1}|+|c_{k+2}d^{k+2}|+\cdots)\\
&<1-e^k+\frac12e^k\\
&=1-\frac12e^k<1.
\end{aligned}
\]
将这样的 $e$ 代入 $h=ed$，得
\[
|f(z_0+ed)|<|f(z_0)|.
\]
这就推出了矛盾. $\square$

\noindent\textbf{推论 5.6.1}\quad 复数域上的一元 $n$ 次多项式恰有 $n$ 个复根 (包括重根).

\noindent\textbf{推论 5.6.2}\quad 复数域上的不可约多项式都是一次多项式.

\noindent\textbf{推论 5.6.3}\quad 复数域上的一元 $n$ 次多项式必可分解为一次因式的乘积.

关于多项式根与系数的关系，我们有如下的 Vieta (韦达) 定理.

\noindent\textbf{定理 5.6.2 (Vieta 定理)}\quad 若数域 $\mathbb{K}$ 上的一元 $n$ 次多项式
\[
f(x)=a_0x^n+a_1x^{n-1}+a_2x^{n-2}+\cdots+a_n
\]
在 $\mathbb{K}$ 中有 $n$ 个根 $x_1,x_2,\cdots,x_n$，则
\[
\sum_{i=1}^n x_i=-\frac{a_1}{a_0},
\]
\[
\sum_{1\leq i<j\leq n}x_ix_j=\frac{a_2}{a_0},
\]
\[
\sum_{1\leq i<j<k\leq n}x_ix_jx_k=-\frac{a_3}{a_0},
\]
\[
\cdots\cdots\cdots
\]
\[
x_1x_2\cdots x_n=(-1)^n\frac{a_n}{a_0}.
\]

\noindent\textbf{证明}\quad $f(x)=a_0(x-x_1)(x-x_2)\cdots(x-x_n)$，将这个式子的右边展开与 $f(x)$ 比较系数即得结论. $\square$

由代数基本定理，一个复数域上的一元 $n$ 次方程必有 $n$ 个根. 如何求这 $n$ 个根，读者已在中学里学到过一元一次及一元二次方程的求根公式. 对于一元三次、一元四次方程也有求根公式，当然要复杂得多. 下面我们介绍一下如何来求解一元三次、一元四次方程.

由于将一个方程式两边乘以非零常数不影响该方程的根，因此不妨设有下列一元三次方程：
\[
f(x)=x^3+ax^2+bx+c=0.
\]
作变换 $x=y-\dfrac13a$，代入上述方程化简后得到一个缺二次项的方程：
\[
y^3+py+q=0.
\]
显然，只要将上面方程的根减去 $\dfrac13a$ 即可得原方程的根，因此我们把问题归结为求
\begin{equation}
f(x)=x^3+px+q=0
\tag{5.6.4}
\end{equation}
这一类方程式的根.

若 $q=0$，则 $x_1=0,\ x_2=\sqrt{-p},\ x_3=-\sqrt{-p}$ 便是方程的根. 若 $p=0$，则 $x_1=\sqrt[3]{-q},\ x_2=\sqrt[3]{-q}\omega,\ x_3=\sqrt[3]{-q}\omega^2$ 就是方程的根，其中
\[
\omega=-\frac12+\frac{\sqrt{3}}{2}i.
\]
因此我们只需讨论 $p\neq 0,\ q\neq 0$ 的情形.

引进新的未知数 $x=u+v$，则
\[
x^3=u^3+v^3+3uv(u+v)=u^3+v^3+3uvx,
\]
或
\[
x^3-3uvx-(u^3+v^3)=0.
\]
与 (5.6.4) 式比较得
\begin{equation}
\left\{
\begin{aligned}
uv&=-\frac13p,\\
u^3+v^3&=-q.
\end{aligned}
\right.
\tag{5.6.5}
\end{equation}
如可求出 (5.6.5) 式中的 $u,v$ 即可求出 $x$，但 (5.6.5) 式又可变为
\begin{equation}
\left\{
\begin{aligned}
u^3v^3&=-\frac1{27}p^3,\\
u^3+v^3&=-q.
\end{aligned}
\right.
\tag{5.6.6}
\end{equation}
由 Vieta 定理可知 $u^3,v^3$ 是下列二次方程的两个根：
\[
y^2+qy-\frac{p^3}{27}=0.
\]
于是
\[
u^3=-\frac q2+\sqrt{\frac{q^2}{4}+\frac{p^3}{27}},\quad
v^3=-\frac q2-\sqrt{\frac{q^2}{4}+\frac{p^3}{27}}.
\]
令
\[
\Delta=\frac{q^2}{4}+\frac{p^3}{27},
\]
注意 $u,v$ 必须适合 (5.6.5) 式，故可得方程 (5.6.4) 的根为
\begin{equation}
\left\{
\begin{aligned}
x_1&=\sqrt[3]{-\frac q2+\sqrt{\Delta}}+\sqrt[3]{-\frac q2-\sqrt{\Delta}},\\
x_2&=\omega\sqrt[3]{-\frac q2+\sqrt{\Delta}}+\omega^2\sqrt[3]{-\frac q2-\sqrt{\Delta}},\\
x_3&=\omega^2\sqrt[3]{-\frac q2+\sqrt{\Delta}}+\omega\sqrt[3]{-\frac q2-\sqrt{\Delta}}.
\end{aligned}
\right.
\tag{5.6.7}
\end{equation}
上式通常称为 Cardano (卡尔达诺) 公式.

现在来考虑四次方程，我们采用与讨论三次方程同样的方法，令 $x=y-\dfrac14a$，可消去方程 $x^4+ax^3+bx^2+cx+d=0$ 的三次项，把求解四次方程的问题归结为求解下面类型的方程：
\begin{equation}
x^4+ax^2+bx+c=0.
\tag{5.6.8}
\end{equation}
引进新的未知数 $u$，在 (5.6.8) 式中加 $ux^2+\dfrac{u^2}{4}$，再减 $ux^2+\dfrac{u^2}{4}$，得
\begin{equation}
x^4+ux^2+\frac{u^2}{4}-\left[(u-a)x^2-bx+\frac{u^2}{4}-c\right]=0.
\tag{5.6.9}
\end{equation}
我们注意到上式中 $x^4+ux^2+\dfrac{u^2}{4}=\left(x^2+\dfrac{u}{2}\right)^2$，如果中括号内是一个完全平方，则 (5.6.9) 式可化为两个二次方程求解. 而中括号是一个完全平方的条件是
\begin{equation}
b^2-4(u-a)\left(\frac{u^2}{4}-c\right)=0.
\tag{5.6.10}
\end{equation}
这是一个 $u$ 的三次方程，称为方程 (5.6.8) 的预解式. 假设 $u$ 已解出，(5.6.9) 式将变成
\[
\left(x^2+\frac u2\right)^2-\left(\sqrt{u-a}x-\frac{b}{2\sqrt{u-a}}\right)^2=0.
\]
分解因式后得到两个二次方程：
\begin{equation}
x^2+\sqrt{u-a}x+\frac u2-\frac{b}{2\sqrt{u-a}}=0,
\tag{5.6.11}
\end{equation}
\begin{equation}
x^2-\sqrt{u-a}x+\frac u2+\frac{b}{2\sqrt{u-a}}=0.
\tag{5.6.12}
\end{equation}
这样便可求出方程 (5.6.8) 的所有根.

这里需要注意的是 (5.6.10) 式是一个三次方程，$u$ 有 3 个根. 如依次将这 3 个根代入 (5.6.11) 式、(5.6.12) 式，岂非得到方程 (5.6.8) 的 12 个根！其实我们只需取 (5.6.10) 式的一个根就可以了. 因为任取其他根所得的结果是完全一样的. 事实上，(5.6.11) 式与 (5.6.12) 式之积就是方程 (5.6.8)，因此方程 (5.6.11)、方程 (5.6.12) 的根总是方程 (5.6.8) 的根，只不过在方程 (5.6.11)、方程 (5.6.12) 中出现的情形不同而已. 比如设 $x_1,x_2,x_3,x_4$ 为方程 (5.6.8) 的 4 个根. 可能 $x_1,x_2$ 为方程 (5.6.11) 的根，这时 $x_3,x_4$ 为方程 (5.6.12) 的根；又可能 $x_3,x_4$ 为方程 (5.6.11) 的根，这时 $x_1,x_2$ 为方程 (5.6.12) 的根，等等. 四次方程的这种解法通常称为 Ferrari (费拉里) 解法.

高于四次以上的方程一般是不能用根式来求解的，这一结论在 19 世纪 30 年代被法国数学家 Galois (伽罗瓦) 证明. 他的证明涉及群、域等抽象代数知识，读者将在抽象代数的课程中学到. 需要注意的是，我们这里说五次及五次以上的方程一般不能用根式求解，但是并不是说不能解. 另外，也并不是所有五次及五次以上的方程都不能用根式求解. 什么时候可以用根式求解，什么时候不能用根式求解，Galois 给出了一个充分必要条件. 读者欲知其详，请参阅有关 Galois 理论的著作，如 [1].

\subsection*{习题 5.6}

1. 设三次方程 $x^3+px^2+qx+r=0$ 的 3 个根成等差数列，求证：
\[
2p^3-9pq+27r=0.
\]

2. 设三次方程 $x^3+px^2+qx+r=0(r\neq 0)$ 的 3 个根成等比数列，求证：
\[
rp^3=q^3.
\]

3. 设多项式 $x^3+3x^2+mx+n$ 的 3 个根成等差数列，多项式 $x^3-(m-2)x^2+(n-3)x+8$ 的 3 个根成等比数列，求 $m$ 和 $n$.

4. 证明：方程 $x^3+px+q=0$ 有重根的充分必要条件是 $4p^3+27q^2=0$.

5. 设方程 $x^3+px^2+qx+r=0$ 的 3 个根都是实数，求证：$p^2\geq 3q$.

6. 设 $f(x)=a_nx^n+a_{n-1}x^{n-1}+\cdots+a_0$ 的 $n$ 个根 $x_1,x_2,\cdots,x_n$ 皆不等于零，求以 $\dfrac1{x_1},\dfrac1{x_2},\cdots,\dfrac1{x_n}$ 为根的多项式.

7. 设 $f(x)$ 是复数域上的多项式，若对任意的实数 $c$，$f(c)$ 总是实数，求证：$f(x)$ 是实系数多项式.

\section*{§5.7\quad 实系数多项式和有理系数多项式}

我们已经知道，复数域上的不可约多项式都是一次的. 实数域上的不可约多项式应该是什么样的呢？

\noindent\textbf{定理 5.7.1}\quad 设
\[
f(x)=a_nx^n+a_{n-1}x^{n-1}+\cdots+a_1x+a_0
\]
是实系数多项式，若复数 $a+bi(b\neq 0)$ 是其根，则 $a-bi$ 也是它的根.

\noindent\textbf{证明}\quad 令 $z=a+bi$，其共轭复数为 $\bar z=a-bi$，则
\[
\begin{aligned}
f(\bar z)&=a_n\bar z^n+a_{n-1}\bar z^{n-1}+\cdots+a_1\bar z+a_0\\
&=\overline{a_nz^n+a_{n-1}z^{n-1}+\cdots+a_1z+a_0}=0.
\end{aligned}
\]
由此即得结论. $\square$

定理 5.7.1 表明，实系数多项式 $f(x)$ 的虚部不为零的复根必成对出现.

\noindent\textbf{推论 5.7.1}\quad 实数域上的不可约多项式为一次多项式或下列二次多项式：
\[
ax^2+bx+c,\quad \text{其中 } b^2-4ac<0.
\]

\noindent\textbf{证明}\quad 一次多项式显然为不可约. 当 $b^2-4ac<0$ 时，$ax^2+bx+c$ 没有实根，故不可约. 反过来，任一高于二次的实系数多项式 $f(x)$ 如有实根，则 $f(x)$ 可约；如有一复根 $a+bi(b\neq 0)$，则 $a-bi$ 也是它的根，从而
\[
(x-(a+bi))(x-(a-bi))=x^2-2ax+(a^2+b^2)
\]
是 $f(x)$ 的因式，故 $f(x)$ 可约. $\square$

\noindent\textbf{推论 5.7.2}\quad 实数域上的多项式 $f(x)$ 必可分解为有限个一次因式及不可约二次因式的乘积.

接下去讨论有理系数多项式. 我们首先证明一个整系数多项式有有理根的必要条件，然后研究有理系数多项式的因式分解.

\noindent\textbf{定理 5.7.2}\quad 设有 $n$ 次整系数多项式
\begin{equation}
f(x)=a_nx^n+a_{n-1}x^{n-1}+\cdots+a_1x+a_0,
\tag{5.7.1}
\end{equation}
则有理数 $\dfrac q p$ 是 $f(x)$ 的根的必要条件是 $p\mid a_n,\ q\mid a_0$，其中 $p,q$ 是互素的整数.

\noindent\textbf{证明}\quad 将 $\dfrac q p$ 代入 (5.7.1) 式得
\[
a_n\left(\frac q p\right)^n+a_{n-1}\left(\frac q p\right)^{n-1}+\cdots+a_1\left(\frac q p\right)+a_0=0,
\]
将上式两边乘以 $p^n$ 得
\[
a_nq^n+a_{n-1}q^{n-1}p+\cdots+a_1qp^{n-1}+a_0p^n=0.
\]
显然必须有 $p\mid a_n,\ q\mid a_0.\quad \square$

定理 5.7.2 给出了整系数多项式有有理根的必要条件. 由于任一有理系数方程均可化为同解的整系数方程 (只需用系数的公分母乘以方程的两边即可)，因此定理 5.7.2 也可用来判别一个有理系数多项式是否有有理根.

\noindent\textbf{例 5.7.1}\quad 证明下列多项式没有有理根：
\[
f(x)=x^5-12x^3+36x+12.
\]

\noindent\textbf{证明}\quad $f(x)$ 如有有理根，只可能为 $\pm1,\pm2,\pm3,\pm4,\pm6,\pm12$. 将上述这些数代入 $f(x)$ 均不为零，因此 $f(x)$ 无有理根. $\square$

现在来讨论有理系数多项式的因式分解. 首先我们要探讨一下有理系数多项式的可约性与整系数多项式可约性的关系.

\noindent\textbf{定义 5.7.1}\quad 设多项式
\[
f(x)=a_nx^n+a_{n-1}x^{n-1}+\cdots+a_1x+a_0
\]
是整系数多项式，若 $a_n,a_{n-1},\cdots,a_1,a_0$ 的最大公约数等于 1，则称 $f(x)$ 为本原多项式.

\noindent\textbf{引理 5.7.1 (Gauss 引理)}\quad 两个本原多项式之积仍是本原多项式.

\noindent\textbf{证明}\quad 设
\[
f(x)=a_nx^n+a_{n-1}x^{n-1}+\cdots+a_1x+a_0,
\]
\[
g(x)=b_mx^m+b_{m-1}x^{m-1}+\cdots+b_1x+b_0
\]
是两个本原多项式. 若
\[
f(x)g(x)=c_{m+n}x^{m+n}+c_{m+n-1}x^{m+n-1}+\cdots+c_1x+c_0
\]
不是本原多项式，则 $c_0,c_1,\cdots,c_{m+n}$ 必有一个公共素因子 $p$. 因为 $f(x)$ 是本原多项式，故 $p$ 不能整除 $f(x)$ 的所有系数，可设 $p\mid a_0,p\mid a_1,\cdots,p\mid a_{i-1}$，但 $p$ 不能整除 $a_i$. 同理，可设 $p\mid b_0,p\mid b_1,\cdots,p\mid b_{j-1}$，但 $p$ 不能整除 $b_j$. 注意到
\[
c_{i+j}=\cdots+a_{i-2}b_{j+2}+a_{i-1}b_{j+1}+a_ib_j+a_{i+1}b_{j-1}+\cdots,
\]
$p$ 可整除 $c_{i+j}$，$p$ 也能整除右式除 $a_ib_j$ 以外的所有项. 但 $p$ 不能整除 $a_i$ 和 $b_j$，故 $p$ 不能整除 $a_ib_j$，引出矛盾. $\square$

\noindent\textbf{定理 5.7.3}\quad 若整系数多项式 $f(x)$ 在有理数域上可约，则它必可分解为两个次数较低的整系数多项式之积.

\noindent\textbf{证明}\quad 假设整系数多项式 $f(x)$ 可以分解为两个次数较低的有理系数多项式之积：
\[
f(x)=g(x)h(x),
\]
$g(x)$ 的各项系数为有理数，必有一个公分母记为 $c$，于是 $g(x)=\dfrac1c(cg(x))$，其中 $cg(x)$ 为整系数多项式. 若把 $cg(x)$ 中所有系数的最大公因数 $d$ 提出来，则
\[
g(x)=\frac dc\left(\frac cdg(x)\right),\quad \frac cdg(x)
\]
是一个本原多项式. 这表明 $g(x)=ag_1(x)$，$a$ 为有理数，$g_1(x)$ 为本原多项式. 同理，$h(x)=bh_1(x)$，其中 $b$ 为有理数，$h_1(x)$ 为本原多项式. 于是我们得到
\[
f(x)=g(x)h(x)=abg_1(x)h_1(x).
\]
由 Gauss 引理知，$g_1(x)h_1(x)$ 是本原多项式. 若 $ab$ 不是一个整数，则 $abg_1(x)h_1(x)$ 将不是整系数多项式，这与 $f(x)$ 是整系数多项式相矛盾. 因此 $ab$ 必须是整数，于是 $f(x)$ 可以分解为两个次数较小的整系数多项式之积. $\square$

我们通常称一个整系数多项式 $f(x)$ 在整数环上可约，若它可以分解为两个次数较低的整系数多项式之积. 定理 5.7.3 告诉我们，整系数多项式 $f(x)$ 若在整数环上不可约，则在有理数域上也不可约. 如何判断一个整系数多项式不可约，我们有如下的 Eisenstein 判别法.

\noindent\textbf{定理 5.7.4 (Eisenstein 判别法)}\quad 设多项式
\[
f(x)=a_nx^n+a_{n-1}x^{n-1}+\cdots+a_1x+a_0
\]
是整系数多项式，$a_n\neq 0,\ n\geq 1,\ p$ 是一个素数. 若 $p\mid a_i(i=0,1,\cdots,n-1)$，但 $p$ 不能整除 $a_n$ 且 $p^2$ 不能整除 $a_0$，则 $f(x)$ 在有理数域上不可约.

\noindent\textbf{证明}\quad 只需证明 $f(x)$ 在整数环上不可约即可. 设 $f(x)$ 可分解为两个次数较低的整系数多项式之积：
\[
f(x)=(b_mx^m+b_{m-1}x^{m-1}+\cdots+b_0)(c_tx^t+c_{t-1}x^{t-1}+\cdots+c_0),
\]
其中 $m+t=n$. 显然 $a_0=b_0c_0,\ a_n=b_mc_t$. 由假设 $p\mid a_0$，故 $p\mid b_0$ 或 $p\mid c_0$. 又 $p^2$ 不能整除 $a_0$，故 $p$ 不能同时整除 $b_0$ 及 $c_0$. 不妨设 $p\mid b_0$ 但 $p$ 不能整除 $c_0$. 又

由假设，$p$ 不能整除 $a_n=b_mc_t$，故 $p$ 既不能整除 $b_m$ 又不能整除 $c_t$. 因此不妨设
\[
p\mid b_0,\ p\mid b_1,\cdots,\ p\mid b_{j-1}
\]
但 $p$ 不能整除 $b_j$，其中 $0<j\leq m<n$. 而
\[
a_j=b_jc_0+b_{j-1}c_1+\cdots+b_0c_j,
\]
根据假设，$p\mid a_j$，又 $p$ 可整除上式右端除 $b_jc_0$ 外的其余项，而不能整除 $b_jc_0$ 这一项，引出矛盾. $\square$

\noindent\textbf{例 5.7.2}\quad 对任意的 $n\geq 1,\ x^n-2$ 在有理数域上不可约.

\noindent\textbf{证明}\quad 由于 2 可整除 $x^n-2$ 中除首项以外的一切系数，又 $2^2=4$ 不能整除 $-2$. 由 Eisenstein 判别法知 $x^n-2$ 不可约. $\square$

这个例子表明，存在任意次的有理数域上的不可约多项式.

\noindent\textbf{例 5.7.3}\quad 若 $p$ 为素数，证明：
\[
f(x)=x^{p-1}+x^{p-2}+\cdots+x+1
\]
在有理数域上不可约.

\noindent\textbf{证明}\quad 作变量代换
\[
x=y+1,
\]
得
\[
f(x)=\frac{x^p-1}{x-1}=\frac{(y+1)^p-1}{y}
=y^{p-1}+C_p^1y^{p-2}+C_p^2y^{p-3}+\cdots+C_p^{p-1}.
\]
注意 $p\mid C_p^i(1\leq i\leq p-1)$，$p$ 不能整除首项系数 1，$p^2$ 不能整除 $C_p^{p-1}=p$，因此，上述关于 $y$ 的多项式在有理数域上不可约，从而 $f(x)$ 在有理数域上也不可约. $\square$

\noindent\textbf{例 5.7.4}\quad 证明：
\[
f(x)=1+x+\frac{x^2}{2!}+\cdots+\frac{x^n}{n!}
\]
当 $n$ 是素数时在有理数域上不可约.

\noindent\textbf{证明}\quad 将 $n!$ 乘以 $f(x)$，只需证明当 $n=p$ 为素数时，在整数环上不可约即可. 而由 Eisenstein 判别法即知 $p!f(x)$ 不可约. $\square$

\subsection*{习题 5.7}

1. 证明：奇数次实系数多项式必有实数根.

2. 设 $f(x)=a_nx^n+a_{n-1}x^{n-1}+\cdots+a_1x+a_0$ 是实系数多项式，求证：

(1) 若 $a_i(0\leq i<n)$ 全是正数或全是负数，则 $f(x)$ 没有非负实根.

(2) 若 $(-1)^ia_i(0\leq i<n)$ 全是正数或全是负数，则 $f(x)$ 没有非正实根.

(3) 若 $a_n>0$ 且 $(-1)^{n-i}a_i>0(0\leq i\leq n-1)$，则 $f(x)$ 没有非正实根；若 $a_n>0$ 且 $(-1)^{n-i}a_i\geq 0(0\leq i\leq n-1)$，则 $f(x)$ 没有负实根.

3. 求证：方程 $x^8+5x^6+4x^4+2x^2+1=0$ 无实数根.

4. 求下列多项式的有理根：

(1) $x^3-6x^2+15x-14$；

(2) $4x^4-7x^2-5x-1$；

(3) $6x^4-25x^3+20x^2-25x+14$；

(4) $x^5-x^4-\dfrac52x^3+2x^2-\dfrac12x-3$.

5. 设 $f(x)$ 是首一整系数多项式，若 $f(0),f(1)$ 都是奇数，求证：$f(x)$ 没有有理根.

6. 证明下列多项式在有理数域上不可约：

(1) $x^p+px+p$ ($p$ 是素数)；

(2) $x^6+x^3+1$；

(3) $x^4-18x^3+12x^2-6x+2$.

7. 设 $p_1,\cdots,p_m$ 是 $m$ 个互不相同的素数，求证：对任意的 $n\geq 1$，下列多项式在有理数域上不可约：
\[
f(x)=x^n-p_1\cdots p_m.
\]

8. 证明：$x^8+1$ 在有理数域上不可约.

9. 写出一个次数最小的首一有理系数多项式，使它有下列根：
\[
1+\sqrt3,\quad 3+\sqrt{2}i.
\]

\section*{§5.8\quad 多元多项式}

设 $\mathbb{K}$ 是数域，$x_1,x_2,\cdots,x_n$ 是 $n$ 个未定元，它们彼此无关. 称 $ax_1^{k_1}x_2^{k_2}\cdots x_n^{k_n}$ 为一个单项式，其中 $a$ 是这个单项式的系数，$k_j$ 为非负整数. 如果 $a\neq 0$，则称该单项式的次数为 $k=k_1+k_2+\cdots+k_n$. 如果两个单项式除相差一个系数外，其余都相同，即每个 $x_i$ 的次数都相同，则称这两个单项式为同类项. 同类项相加，可将它们的系数相加. 比如
\[
2x_1^2x_2x_3^4+3x_1^2x_2x_3^4=5x_1^2x_2x_3^4.
\]
不是同类项相加时不能把系数相加. 有限个单项式的和称为一个 $n$ 元多项式，它的一般形式为
\[
f(x_1,x_2,\cdots,x_n)=\sum a_{i_1i_2\cdots i_n}x_1^{i_1}x_2^{i_2}\cdots x_n^{i_n}.
\]
下面我们总假设在一个文字多项式中的各单项式彼此不同，即同类项已合并.

对一个 $n$ 元多项式，它的系数非零的单项式的最大次数称为这个多项式的次数. 比如下列三元多项式的次数为 5：
\[
x_1^2x_2x_3^2+2x_1x_2x_3-3x_1x_2+1.
\]
两个 $n$ 元多项式相等当且仅当它们同类项的系数全都相等.

两个 $n$ 元多项式相加，即将它们同类项的系数相加. 例如，若
\[
f(x_1,x_2,x_3)=x_1^3+2x_1x_3-3x_1x_2x_3,
\]
\[
g(x_1,x_2,x_3)=-x_2^2-2x_1x_3+5x_1x_2x_3,
\]
则
\[
f(x_1,x_2,x_3)+g(x_1,x_2,x_3)=x_1^3-x_2^2+2x_1x_2x_3.
\]
两个单项式 $ax_1^{i_1}x_2^{i_2}\cdots x_n^{i_n}$，$bx_1^{j_1}x_2^{j_2}\cdots x_n^{j_n}$ 相乘，其积为
\[
abx_1^{i_1+j_1}x_2^{i_2+j_2}\cdots x_n^{i_n+j_n}.
\]
两个多项式相乘按分配律可化为各单项式乘积之和. 如
\[
(x_1^3+x_1x_2x_3-2x_3)(x_1^2-x_1x_2-x_3^2)
\]
\[
=x_1^5-x_1^4x_2-x_1^3x_3^2+x_1^3x_2x_3-x_1^2x_2^2x_3-x_1x_2x_3^3-2x_1^2x_3+2x_1x_2x_3+2x_3^3.
\]
相乘后如出现同类项应予以合并. 多元多项式不像一元多项式那样可按次数的大小降幂或升幂排列. 比如
\[
x_1^2x_2+x_1x_2x_3,
\]
其中两项都是三次式. 为了给它们排次序，我们常采用“字典排列法”：$n$ 个未定元按自然足标为序排列，即为 $x_1,x_2,\cdots,x_n$，若有两个非零单项式：
\[
ax_1^{i_1}x_2^{i_2}\cdots x_n^{i_n},\quad bx_1^{j_1}x_2^{j_2}\cdots x_n^{j_n},
\]
若 $i_1=j_1,\ i_2=j_2,\cdots,\ i_{k-1}=j_{k-1}$，但 $i_k>j_k$，则规定 $ax_1^{i_1}x_2^{i_2}\cdots x_n^{i_n}$ 先于 $bx_1^{j_1}x_2^{j_2}\cdots x_n^{j_n}$. 按这样排列，任一多项式都只有唯一的方法把它的各单项式排序. 这时要注意第一项即首项未必是最高次项，末项也未必是次数最低的项. 例如多项式：
\[
x_1^2x_2+x_1^3x_3+3x_1x_2x_3-4x_2x_3^5
\]
按字典排列应为
\[
x_1^3x_3+x_1^2x_2+3x_1x_2x_3-4x_2x_3^5.
\]

\noindent\textbf{引理 5.8.1}\quad 若 $f(x_1,x_2,\cdots,x_n)$ 及 $g(x_1,x_2,\cdots,x_n)$ 都是 $\mathbb{K}$ 上非零的 $n$ 元多项式，则按字典排列法排列后乘积的首项等于 $f$ 的首项与 $g$ 的首项之积.

\noindent\textbf{证明}\quad 设 $ax_1^{i_1}x_2^{i_2}\cdots x_n^{i_n}$ 和 $bx_1^{j_1}x_2^{j_2}\cdots x_n^{j_n}$ 分别是 $f$ 和 $g$ 的首项 (按字典排列法)，它们的乘积为 $abx_1^{i_1+j_1}x_2^{i_2+j_2}\cdots x_n^{i_n+j_n}$. 其他任意两个单项式 $cx_1^{k_1}x_2^{k_2}\cdots x_n^{k_n}$ 和 $dx_1^{r_1}x_2^{r_2}\cdots x_n^{r_n}$ 之积为 $cdx_1^{k_1+r_1}x_2^{k_2+r_2}\cdots x_n^{k_n+r_n}$. 设 $i_1=k_1,\cdots,i_{t-1}=k_{t-1}$，$i_t>k_t$；$j_1=r_1,\cdots,j_{s-1}=r_{s-1}$，$j_s>r_s$. 不妨设 $t\leq s$，显然
\[
i_1+j_1=k_1+r_1,\ \cdots,\ i_{t-1}+j_{t-1}=k_{t-1}+r_{t-1},\ i_t+j_t>k_t+r_t.
\]
因此 $abx_1^{i_1+j_1}x_2^{i_2+j_2}\cdots x_n^{i_n+j_n}$ 先于 $cdx_1^{k_1+r_1}x_2^{k_2+r_2}\cdots x_n^{k_n+r_n}$.

同理可证明：$abx_1^{i_1+j_1}x_2^{i_2+j_2}\cdots x_n^{i_n+j_n}$ 先于 $adx_1^{i_1+r_1}x_2^{i_2+r_2}\cdots x_n^{i_n+r_n}$ 和 $cbx_1^{k_1+j_1}x_2^{k_2+j_2}\cdots x_n^{k_n+j_n}$. 因此它确是 $fg$ 的首项. $\square$

\noindent\textbf{命题 5.8.1}\quad 若 $f(x_1,x_2,\cdots,x_n)\neq 0,\ g(x_1,x_2,\cdots,x_n)\neq 0$，则
\[
f(x_1,x_2,\cdots,x_n)g(x_1,x_2,\cdots,x_n)\neq 0.
\]

\noindent\textbf{证明}\quad 由 $f$ 和 $g$ 的首项不为零可知 $fg$ 的首项不为零，于是 $fg\neq 0$. $\square$

\noindent\textbf{推论 5.8.1}\quad 若 $h(x_1,x_2,\cdots,x_n)\neq 0$，且
\[
f(x_1,x_2,\cdots,x_n)h(x_1,x_2,\cdots,x_n)=g(x_1,x_2,\cdots,x_n)h(x_1,x_2,\cdots,x_n),
\]
则
\[
f(x_1,x_2,\cdots,x_n)=g(x_1,x_2,\cdots,x_n).
\]

多元多项式除了“字典排列法”外，还有“齐次排列法”. 我们先介绍齐次多项式的概念. 若一个多项式 $f(x_1,x_2,\cdots,x_n)$ 的每个单项式都是 $k$ 次式，则称之为 $k$ 次齐次多项式或 $k$ 次型. 如 $a_1x_1+a_2x_2+a_3x_3$ 是三元一次型. $x_1^2+x_1x_2-3x_2^2$ 为二元二次型. $x_1^4-4x_2x_3x_4^2+5x_1x_2x_3x_4$ 为四元四次型.

两个次数相同的齐次多项式之和若不为零，则必仍是同次齐次多项式. 任意两个齐次多项式之积仍为齐次多项式.

任一 $n$ 元多项式均可表示为若干个齐次多项式之和，这只需要将各次数相等的项放在一起即可. 如
\[
2x_1^3-3x_1x_2+4x_4^3-x_1x_2x_3+x_4^2=(2x_1^3+4x_4^3-x_1x_2x_3)+(x_4^2-3x_1x_2),
\]
其中 $2x_1^3+4x_4^3-x_1x_2x_3$ 为三次型，$x_4^2-3x_1x_2$ 为二次型.

多元多项式与多元多项式函数之间的关系与一元的情形类似.

\noindent\textbf{引理 5.8.2}\quad 设 $f(x_1,x_2,\cdots,x_n)$ 是 $\mathbb{K}$ 上非零的 $n$ 元多项式，则必存在 $\mathbb{K}$ 中的数 $a_1,a_2,\cdots,a_n$，使 $f(a_1,a_2,\cdots,a_n)\neq 0$.

\noindent\textbf{证明}\quad 对未定元个数 $n$ 用数学归纳法. 当 $n=1$ 时，多项式 $f(x)$ 只有有限个零点，故总有 $a\in\mathbb{K}$ 使 $f(a)\neq 0$. 现设对有 $n-1$ 个未定元的多项式结论成立. 将 $f(x_1,x_2,\cdots,x_n)$ 写成未定元 $x_n$ 的多项式：
\[
f(x_1,x_2,\cdots,x_n)=b_0+b_1x_n+\cdots+b_mx_n^m,
\]
其中 $b_i=b_i(x_1,x_2,\cdots,x_{n-1})$ 是 $n-1$ 元多项式. 因为 $f(x_1,x_2,\cdots,x_n)\neq 0$，故可设 $b_m\neq 0$. 由归纳假设，存在 $a_1,\cdots,a_{n-1}\in\mathbb{K}$，使
\[
b_m(a_1,\cdots,a_{n-1})\neq 0.
\]
因而
\[
f(a_1,\cdots,a_{n-1},x_n)=b_0(a_1,\cdots,a_{n-1})+b_1(a_1,\cdots,a_{n-1})x_n+\cdots+b_m(a_1,\cdots,a_{n-1})x_n^m
\]
是一个非零的以 $x_n$ 为未定元的一元多项式，故存在 $a_n\in\mathbb{K}$，使
\[
f(a_1,a_2,\cdots,a_n)\neq 0.\quad \square
\]

\noindent\textbf{命题 5.8.2}\quad 数域 $\mathbb{K}$ 上的两个 $n$ 元多项式 $f(x_1,x_2,\cdots,x_n)$ 与 $g(x_1,x_2,\cdots,x_n)$ 相等的充分必要条件是对一切 $a_1,a_2,\cdots,a_n\in\mathbb{K}$，均有
\[
f(a_1,a_2,\cdots,a_n)=g(a_1,a_2,\cdots,a_n).
\]

\noindent\textbf{证明}\quad 只需证明充分性. 作
\[
h(x_1,x_2,\cdots,x_n)=f(x_1,x_2,\cdots,x_n)-g(x_1,x_2,\cdots,x_n).
\]
若 $h(x_1,x_2,\cdots,x_n)\neq 0$，则必有 $a_1,a_2,\cdots,a_n\in\mathbb{K}$，使 $h(a_1,a_2,\cdots,a_n)\neq 0$，这与假设矛盾. $\square$

\subsection*{习题 5.8}

1. 设 $f(x_1,x_2,\cdots,x_n),g(x_1,x_2,\cdots,x_n)$ 是数域 $\mathbb{K}$ 上的两个多项式且 $g(x_1,x_2,\cdots,x_n)$ 是非零多项式. 如果对一切使 $g(a_1,a_2,\cdots,a_n)\neq 0$ 的 $a_1,a_2,\cdots,a_n\in\mathbb{K}$，均有 $f(a_1,a_2,\cdots,a_n)=0$，求证：$f(x_1,x_2,\cdots,x_n)=0$.

2. 设 $f(x_1,x_2,\cdots,x_n),g(x_1,x_2,\cdots,x_n)$ 是数域 $\mathbb{K}$ 上的两个多项式且 $fg$ 是 $n$ 元齐次多项式，求证：$f$ 与 $g$ 都是齐次多项式.

3. 设 $f(x_1,x_2,\cdots,x_n)$ 是 $n$ 元多项式，若存在 $n$ 元多项式 $g(x_1,x_2,\cdots,x_n)$，使
\[
f(x_1,x_2,\cdots,x_n)g(x_1,x_2,\cdots,x_n)=1,
\]
求证：$f(x_1,x_2,\cdots,x_n)=c$，其中 $c$ 是 $\mathbb{K}$ 中的非零元.

\section*{§5.9\quad 对称多项式}

\noindent\textbf{定义 5.9.1}\quad 设 $f(x_1,x_2,\cdots,x_n)$ 是数域 $\mathbb{K}$ 上的 $n$ 元多项式，若对任意的 $1\leq i<j\leq n$，均有
\[
f(x_1,\cdots,x_i,\cdots,x_j,\cdots,x_n)=f(x_1,\cdots,x_j,\cdots,x_i,\cdots,x_n),
\]
则称 $f(x_1,x_2,\cdots,x_n)$ 是数域 $\mathbb{K}$ 上的 $n$ 元对称多项式.

例如，三元多项式 $x_1^2+x_2^2+x_3^2$，将 $x_1$ 与 $x_2$ 对换有
\[
x_1^2+x_2^2+x_3^2=x_2^2+x_1^2+x_3^2.
\]
又将 $x_1$ 与 $x_3$ 对换及将 $x_2$ 与 $x_3$ 对换都得到同一个多项式. 因此 $x_1^2+x_2^2+x_3^2$ 是对称多项式. 多项式 $x_1^2-x_1x_2$ 与 $x_2^2-x_1x_2$ 不相等，故 $x_1^2-x_1x_2$ 不是对称多项式.

设 $(k_1,k_2,\cdots,k_n)$ 是数组 $(1,2,\cdots,n)$ 的一个全排列. 若 $f(x_1,x_2,\cdots,x_n)$ 是一个对称多项式，则不难看出：
\[
f(x_{k_1},x_{k_2},\cdots,x_{k_n})=f(x_1,x_2,\cdots,x_n).
\]
我们称 $x_1\to x_{k_1},\ x_2\to x_{k_2},\cdots,\ x_n\to x_{k_n}$ 是未定元的一个置换. 因此对称多项式在未定元的任一置换下不变.

容易看出对称多项式之和仍是对称多项式，对称多项式之积也是对称多项式. 因此对称多项式的多项式还是对称多项式. 这句话的意思是说：若 $f_1,f_2,\cdots,f_m$ 是 $m$ 个 $n$ 元对称多项式，$g(y_1,y_2,\cdots,y_m)$ 是 $m$ 元多项式，则将 $f_1,f_2,\cdots,f_m$ 代替 $y_1,y_2,\cdots,y_m$ 后得到的多项式
\[
h(x_1,x_2,\cdots,x_n)=g(f_1,f_2,\cdots,f_m)
\]
仍是一个 $n$ 元对称多项式.

在对称多项式中，有一类基本的多项式，称为初等对称多项式. 它们是这样定义的：
\[
\sigma_1=x_1+x_2+\cdots+x_n=\sum_{i=1}^n x_i,
\]
\[
\sigma_2=x_1x_2+x_1x_3+\cdots+x_{n-1}x_n=\sum_{1\leq i<j\leq n}x_ix_j,
\]
\[
\cdots\cdots\cdots
\]
\[
\sigma_n=x_1x_2\cdots x_n.
\]
这 $n$ 个多项式称为 $n$ 元初等对称多项式. 初等对称多项式之所以重要，是因为我们有下列定理.

\noindent\textbf{定理 5.9.1}\quad 设 $f(x_1,x_2,\cdots,x_n)$ 是数域 $\mathbb{K}$ 上的对称多项式，则必存在 $\mathbb{K}$ 上唯一的一个多项式 $g(y_1,y_2,\cdots,y_n)$，使
\[
f(x_1,x_2,\cdots,x_n)=g(\sigma_1,\sigma_2,\cdots,\sigma_n).
\]

\noindent\textbf{证明}\quad 存在性：设 $f(x_1,x_2,\cdots,x_n)$ 按字典排列法的首项为
\[
ax_1^{i_1}x_2^{i_2}\cdots x_n^{i_n},\quad a\neq 0,
\]
则必有 $i_1\geq i_2\geq \cdots\geq i_n$. 事实上 $f(x_1,x_2,\cdots,x_n)$ 是对称多项式，若 $i_k<i_{k+1}$，则将 $x_k$ 与 $x_{k+1}$ 对换得到 $ax_1^{i_1}\cdots x_k^{i_{k+1}}x_{k+1}^{i_k}\cdots x_n^{i_n}$，这一项也是 $f(x_1,x_2,\cdots,x_n)$ 中的项，但它在 $ax_1^{i_1}x_2^{i_2}\cdots x_n^{i_n}$ 之前，此与首项的假设相矛盾.

作多项式
\[
g_1(x_1,x_2,\cdots,x_n)=a\sigma_1^{i_1-i_2}\sigma_2^{i_2-i_3}\cdots\sigma_{n-1}^{i_{n-1}-i_n}\sigma_n^{i_n}.
\]
显然 $g_1$ 是 $x_1,x_2,\cdots,x_n$ 的对称多项式，且 $g_1$ 的首项为
\[
ax_1^{i_1-i_2}(x_1x_2)^{i_2-i_3}\cdots(x_1\cdots x_n)^{i_n}
=ax_1^{i_1}x_2^{i_2}\cdots x_n^{i_n}.
\]
因此 $g_1$ 与 $f$ 的首项相同，于是 $f_1=f-g_1$ 是一个对称多项式，其首项后于 $f$ 的首项. 对 $f_1$ 重复上述做法，不断做下去，便得到一列对称多项式：
\[
f_0=f,\ f_1=f_0-g_1,\ f_2=f_1-g_2,\ \cdots,
\]
每个 $f_i$ 的首项都后于 $f_{i-1}$ 的首项. 设 $bx_1^{k_1}x_2^{k_2}\cdots x_n^{k_n}$ 是某个 $f_i(i\geq 1)$ 的首项，因它后于 $f$ 的首项，故有
\[
i_1\geq k_1\geq k_2\geq \cdots\geq k_n\geq 0.
\]
这样的 $k_1,k_2,\cdots,k_n$ 只有有限个，故多项式 $f_i$ 不能无限地构造下去，即存在某个正整数 $s$，使 $f_s=0$. 于是
\[
f=g_1+f_1=g_1+g_2+f_2=\cdots=g_1+g_2+\cdots+g_s.
\]
由于每个 $g_i$ 都可表示为 $\sigma_1,\sigma_2,\cdots,\sigma_n$ 的多项式，故 $f$ 也可表示为 $\sigma_1,\sigma_2,\cdots,\sigma_n$ 的多项式.

唯一性：设 $g(y_1,\cdots,y_n)$ 及 $h(y_1,\cdots,y_n)$ 都是 $\mathbb{K}$ 上的 $n$ 元多项式，且
\[
f(x_1,\cdots,x_n)=g(\sigma_1,\cdots,\sigma_n)=h(\sigma_1,\cdots,\sigma_n).
\]
令
\[
\varphi(y_1,\cdots,y_n)=g(y_1,\cdots,y_n)-h(y_1,\cdots,y_n),
\]
由假设，有
\[
\varphi(\sigma_1,\cdots,\sigma_n)=0.
\]
我们要证明多项式
\[
\varphi(y_1,\cdots,y_n)=0.
\]
假设 $\varphi(y_1,\cdots,y_n)\neq 0$，不妨设
\[
\varphi(y_1,y_2,\cdots,y_n)=cy_1^{k_1}y_2^{k_2}\cdots y_n^{k_n}+dy_1^{j_1}y_2^{j_2}\cdots y_n^{j_n}+\cdots,
\]
其中 $c,d,\cdots$ 均不为零且假设各单项式彼此不是同类项. 在 $\varphi(\sigma_1,\sigma_2,\cdots,\sigma_n)$ 中，
\[
c\sigma_1^{k_1}\sigma_2^{k_2}\cdots\sigma_n^{k_n}
=cx_1^{k_1}(x_1x_2)^{k_2}\cdots(x_1x_2\cdots x_n)^{k_n}+\cdots
\]
\[
=cx_1^{k_1+k_2+\cdots+k_n}x_2^{k_2+k_3+\cdots+k_n}\cdots x_n^{k_n}+\cdots,
\]
\[
d\sigma_1^{j_1}\sigma_2^{j_2}\cdots\sigma_n^{j_n}
=dx_1^{j_1}(x_1x_2)^{j_2}\cdots(x_1x_2\cdots x_n)^{j_n}+\cdots
\]
\[
=dx_1^{j_1+j_2+\cdots+j_n}x_2^{j_2+j_3+\cdots+j_n}\cdots x_n^{j_n}+\cdots.
\]
因此 $c\sigma_1^{k_1}\sigma_2^{k_2}\cdots\sigma_n^{k_n}$ 与 $d\sigma_1^{j_1}\sigma_2^{j_2}\cdots\sigma_n^{j_n}$ 等化成 $x_1,x_2,\cdots,x_n$ 的多项式后其首项都不是同类项，从而 $\varphi(\sigma_1,\sigma_2,\cdots,\sigma_n)$ 的首项一定是这些首项中排在最前的那一个. 特别，
\[
\varphi(\sigma_1,\sigma_2,\cdots,\sigma_n)\neq 0,
\]
引出矛盾. $\square$

上述定理通常被称为对称多项式基本定理. 这个定理的存在性证明是构造性的，可用来求对称多项式的初等对称多项式表示.

\noindent\textbf{例 5.9.1}\quad 将对称多项式
\[
f(x_1,x_2,x_3)=x_1^2x_2+x_1^2x_3+x_1x_2^2+x_1x_3^2+x_2^2x_3+x_2x_3^2
\]
表示为 $\sigma_1,\sigma_2,\sigma_3$ 的多项式.

\noindent\textbf{解}\quad 这时首项为 $x_1^2x_2$，作
\[
\sigma_1^{2-1}\sigma_2^1=\sigma_1\sigma_2=(x_1+x_2+x_3)(x_1x_2+x_1x_3+x_2x_3)
\]
\[
=x_1^2x_2+x_1^2x_3+x_1x_2^2+x_1x_3^2+x_2^2x_3+x_2x_3^2+3x_1x_2x_3.
\]
因此
\[
f(x_1,x_2,x_3)=\sigma_1\sigma_2-3\sigma_3.
\]

这种做法当多项式次数较高时计算可能相当繁琐. 下面我们通过举例介绍“待定系数法”.

\noindent\textbf{例 5.9.2}\quad 试将对称多项式
\[
f(x_1,x_2,x_3)=(x_1^2+x_2^2)(x_1^2+x_3^2)(x_2^2+x_3^2)
\]
表示为初等对称多项式的多项式.

\noindent\textbf{解}\quad 注意到 $f$ 是一个齐次多项式，次数等于 6. 又 $f$ 的首项是 $x_1^4x_2^2$，它的指数组为 $(4,2,0)$. 从定理 5.9.1 的证明中可看出 $f_i$ 首项的指数组只可能是 $(4,1,1),(3,3,0),(3,2,1),(2,2,2)$，相对应的 $f_i$ 的首项为
\[
\sigma_1^{4-1}\sigma_2^{1-1}\sigma_3^1=\sigma_1^3\sigma_3,\quad
\sigma_1^{3-3}\sigma_2^{3-0}\sigma_3^0=\sigma_2^3,\quad
\sigma_1^{3-2}\sigma_2^{2-1}\sigma_3^1=\sigma_1\sigma_2\sigma_3,\quad
\sigma_1^{2-2}\sigma_2^{2-2}\sigma_3^2=\sigma_3^2,
\]
因此可设
\[
f=\sigma_1^2\sigma_2^2+a\sigma_1^3\sigma_3+b\sigma_2^3+c\sigma_1\sigma_2\sigma_3+d\sigma_3^2.
\]
取 $x_1,x_2,x_3$ 的一些特殊值便得到关于 $a,b,c,d$ 的线性方程组. 不难解得
\[
a=-2,\ b=-2,\ c=4,\ d=-1.
\]
因此
\[
f=\sigma_1^2\sigma_2^2-2\sigma_1^3\sigma_3-2\sigma_2^3+4\sigma_1\sigma_2\sigma_3-\sigma_3^2.
\]

上述例子中 $f$ 是一个齐次多项式. 若 $f$ 不是齐次多项式，则可把 $f$ 分解为齐次多项式之和. 显然每个齐次多项式仍是对称多项式，可用初等对称多项式来表示. 这样便可得到 $f$ 的表示.

下面我们证明著名的 Newton (牛顿) 公式. 令
\[
s_k=x_1^k+x_2^k+\cdots+x_n^k\ (k\geq 1);\quad s_0=n.
\]

\noindent\textbf{引理 5.9.1}\quad 设
\[
f(x)=(x-x_1)(x-x_2)\cdots(x-x_n)
\]
\[
=x^n-\sigma_1x^{n-1}+\sigma_2x^{n-2}+\cdots+(-1)^n\sigma_n,
\]
则
\[
x^{k+1}f'(x)=(s_0x^k+s_1x^{k-1}+\cdots+s_k)f(x)+g(x),
\]
其中 $g(x)$ 作为 $x$ 的多项式次数小于 $n$.

\noindent\textbf{证明}\quad 容易看出
\[
f'(x)=\sum_{i=1}^n\frac{f(x)}{x-x_i}.
\]
因此
\[
x^{k+1}f'(x)=\sum_{i=1}^n\frac{x^{k+1}}{x-x_i}f(x)
\]
\[
=\sum_{i=1}^n\frac{x^{k+1}-x_i^{k+1}+x_i^{k+1}}{x-x_i}f(x)
\]
\[
=f(x)\sum_{i=1}^n\frac{x^{k+1}-x_i^{k+1}}{x-x_i}+g(x),
\]
其中
\[
g(x)=\sum_{i=1}^n\frac{x_i^{k+1}f(x)}{x-x_i}
\]
作为 $x$ 的多项式其次数小于 $n$. 又
\[
\sum_{i=1}^n\frac{x^{k+1}-x_i^{k+1}}{x-x_i}
=\sum_{i=1}^n(x^k+x_ix^{k-1}+\cdots+x_i^k)
\]
\[
=nx^k+(x_1+\cdots+x_n)x^{k-1}+\cdots+(x_1^k+\cdots+x_n^k)
\]
\[
=s_0x^k+s_1x^{k-1}+\cdots+s_k,
\]
于是
\[
x^{k+1}f'(x)=(s_0x^k+s_1x^{k-1}+\cdots+s_k)f(x)+g(x).\quad \square
\]

\noindent\textbf{命题 5.9.1 (Newton 公式)}\quad 记号同上，若 $k\leq n-1$，则
\[
s_k-s_{k-1}\sigma_1+s_{k-2}\sigma_2-\cdots+(-1)^{k-1}s_1\sigma_{k-1}+(-1)^kk\sigma_k=0;
\]
若 $k\geq n$，则
\[
s_k-s_{k-1}\sigma_1+s_{k-2}\sigma_2-\cdots+(-1)^ns_{k-n}\sigma_n=0.
\]

\noindent\textbf{证明}\quad 对 $f(x)=x^n-\sigma_1x^{n-1}+\sigma_2x^{n-2}+\cdots+(-1)^n\sigma_n$ 求导并乘以 $x^{k+1}$ 得到
\[
x^{k+1}f'(x)=nx^{n+k}-(n-1)\sigma_1x^{n+k-1}+(n-2)\sigma_2x^{n+k-2}-\cdots+(-1)^{n-1}\sigma_{n-1}x^{k+1}.
\]
由引理，有
\[
x^{k+1}f'(x)=(s_0x^k+s_1x^{k-1}+\cdots+s_k)f(x)+g(x)
\]
\[
=(s_0x^k+s_1x^{k-1}+\cdots+s_k)(x^n-\sigma_1x^{n-1}
\]
\[
+\sigma_2x^{n-2}-\cdots+(-1)^n\sigma_n)+g(x).
\]
比较 $x^n$ 的系数即知，当 $k\leq n-1$ 时，有
\[
s_k-s_{k-1}\sigma_1+s_{k-2}\sigma_2-\cdots+(-1)^{k-1}s_1\sigma_{k-1}+(-1)^kk\sigma_k=0;
\]
若 $k\geq n$，则
\[
s_k-s_{k-1}\sigma_1+s_{k-2}\sigma_2-\cdots+(-1)^ns_{k-n}\sigma_n=0.\quad \square
\]

\subsection*{习题 5.9}

1. 求下列对称多项式的初等对称多项式表示：

(1) $(x_1+x_2+x_3)^3-(x_2+x_3-x_1)^3-(x_1+x_3-x_2)^3-(x_1+x_2-x_3)^3$；

(2) $x_1^3+x_2^3+x_3^3-3x_1x_2x_3$.

2. 已知 $x^3+px^2+qx+r=0$ 的 3 个根为 $x_1,x_2,x_3$，求一个三次方程，其根为 $x_1^3,x_2^3,x_3^3$.

3. 求证：$f(x)=x^3+px^2+qx+r$ 的某一根平方等于其他两根平方和的充分必要条件是
\[
p^4(p^2-2q)=2(p^3-2pq+2r)^2.
\]

4. 利用 Newton 公式求出 $s_4$ 的初等对称多项式表示.

5. 求解下列方程组：
\[
\begin{cases}
x_1+x_2+x_3+x_4=4,\\
x_1^2+x_2^2+x_3^2+x_4^2=4,\\
x_1^3+x_2^3+x_3^3+x_4^3=4,\\
x_1^4+x_2^4+x_3^4+x_4^4=4.
\end{cases}
\]

6. 设 $1\leq k\leq n$，求证：
\[
(1)\quad \sigma_k=\frac1{k!}
\begin{vmatrix}
s_1&1&0&\cdots&0\\
s_2&s_1&2&\cdots&0\\
\vdots&\vdots&\vdots&&\vdots\\
s_{k-1}&s_{k-2}&s_{k-3}&\cdots&k-1\\
s_k&s_{k-1}&s_{k-2}&\cdots&s_1
\end{vmatrix};
\]
\[
(2)\quad s_k=
\begin{vmatrix}
\sigma_1&1&0&\cdots&0\\
2\sigma_2&\sigma_1&1&\cdots&0\\
\vdots&\vdots&\vdots&&\vdots\\
(k-1)\sigma_{k-1}&\sigma_{k-2}&\sigma_{k-3}&\cdots&1\\
k\sigma_k&\sigma_{k-1}&\sigma_{k-2}&\cdots&\sigma_1
\end{vmatrix}.
\]

\section*{§5.10\quad 结式和判别式}

设 $f(x),g(x)$ 是数域 $\mathbb{K}$ 上的一元多项式，现在我们来讨论它们何时有公共根 (简称公根). 公根问题实际上等价于公因式问题，但是现在我们要从另外一个角度来探讨这个问题.

\noindent\textbf{引理 5.10.1}\quad 设 $d(x)$ 是 $f(x)$ 与 $g(x)$ 的最大公因式，则 $d(x)\neq 1$ 的充分必要条件是存在 $\mathbb{K}$ 上的非零多项式 $u(x),v(x)$，使
\begin{equation}
f(x)u(x)=g(x)v(x),
\tag{5.10.1}
\end{equation}
且 $\deg u(x)<\deg g(x),\ \deg v(x)<\deg f(x)$.

\noindent\textbf{证明}\quad 若 $d(x)\neq 1$，令 $f(x)=d(x)v(x),\ g(x)=d(x)u(x)$，则
\[
f(x)u(x)=d(x)v(x)u(x)=g(x)v(x),
\]
且 $\deg u(x)<\deg g(x),\ \deg v(x)<\deg f(x)$.

反过来，若 $d(x)=1$，则由 (5.10.1) 式知 $f(x)\mid g(x)v(x)$，由于 $f(x)$ 与 $g(x)$ 互素，故 $f(x)\mid v(x)$，这与 $\deg f(x)>\deg v(x)$ 矛盾. $\square$

现设
\[
\begin{aligned}
f(x)&=a_0x^n+a_1x^{n-1}+\cdots+a_n,\\
g(x)&=b_0x^m+b_1x^{m-1}+\cdots+b_m,\\
u(x)&=x_0x^{m-1}+x_1x^{m-2}+\cdots+x_{m-1},\\
v(x)&=y_0x^{n-1}+y_1x^{n-2}+\cdots+y_{n-1},
\end{aligned}
\]
其中 $x_0,\cdots,x_{m-1};y_0,\cdots,y_{n-1}$ 为待定未知数. 将上面 4 个式子代入 (5.10.1) 式，比较系数得
\[
\begin{cases}
a_0x_0=b_0y_0,\\
a_1x_0+a_0x_1=b_1y_0+b_0y_1,\\
a_2x_0+a_1x_1+a_0x_2=b_2y_0+b_1y_1+b_0y_2,\\
\cdots\cdots\cdots\\
a_nx_{m-3}+a_{n-1}x_{m-2}+a_{n-2}x_{m-1}=b_my_{n-3}+b_{m-1}y_{n-2}+b_{m-2}y_{n-1},\\
a_nx_{m-2}+a_{n-1}x_{m-1}=b_my_{n-2}+b_{m-1}y_{n-1},\\
a_nx_{m-1}=b_my_{n-1}.
\end{cases}
\]
我们把上述 $m+n$ 个等式看成是 $m+n$ 个未知数 $x_0,x_1,\cdots,x_{m-1};y_0,y_1,\cdots,y_{n-1}$ 的线性方程组. 不难算出这个线性方程组系数矩阵的转置为
\[
\begin{pmatrix}
a_0&a_1&a_2&\cdots&\cdots&a_n&0&\cdots&0\\
0&a_0&a_1&\cdots&\cdots&a_{n-1}&a_n&\cdots&0\\
0&0&a_0&\cdots&\cdots&a_{n-2}&a_{n-1}&\cdots&0\\
\vdots&\vdots&\vdots&\vdots&\vdots&\vdots&\vdots&\vdots&\vdots\\
0&0&\cdots&0&a_0&\cdots&\cdots&\cdots&a_n\\
-b_0&-b_1&-b_2&\cdots&\cdots&\cdots&-b_m&\cdots&0\\
0&-b_0&-b_1&\cdots&\cdots&\cdots&-b_{m-1}&-b_m&\cdots\\
\vdots&\vdots&\vdots&\vdots&\vdots&\vdots&\vdots&\vdots&\vdots\\
0&\cdots&0&-b_0&-b_1&\cdots&\cdots&\cdots&-b_m
\end{pmatrix}.
\]
若上述 $m+n$ 阶方阵的行列式不等于零，则 $x_i,y_i$ 都只能全为零，这时 $f(x),g(x)$ 互素，即没有公根. 反之，若上述方阵的行列式等于零，则 $f(x),g(x)$ 有非常数公因式，即有公根.

\noindent\textbf{定义 5.10.1}\quad 设
\[
f(x)=a_0x^n+a_1x^{n-1}+\cdots+a_n,
\]
\[
g(x)=b_0x^m+b_1x^{m-1}+\cdots+b_m.
\]
定义下列 $m+n$ 阶行列式：
\[
R(f,g)=
\begin{vmatrix}
a_0&a_1&a_2&\cdots&\cdots&a_n&0&\cdots&0\\
0&a_0&a_1&\cdots&\cdots&a_{n-1}&a_n&\cdots&0\\
0&0&a_0&\cdots&\cdots&a_{n-2}&a_{n-1}&\cdots&0\\
\vdots&\vdots&\vdots&\vdots&\vdots&\vdots&\vdots&\vdots&\vdots\\
0&0&\cdots&0&a_0&\cdots&\cdots&\cdots&a_n\\
b_0&b_1&b_2&\cdots&\cdots&\cdots&b_m&\cdots&0\\
0&b_0&b_1&\cdots&\cdots&\cdots&b_{m-1}&b_m&\cdots\\
\vdots&\vdots&\vdots&\vdots&\vdots&\vdots&\vdots&\vdots&\vdots\\
0&\cdots&0&b_0&b_1&\cdots&\cdots&\cdots&b_m
\end{vmatrix}
\]
为 $f(x)$ 与 $g(x)$ 的结式或称为 Sylvester 行列式.

显然我们可以有下列判断两个多项式存在公根的定理.

\noindent\textbf{定理 5.10.1}\quad 多项式 $f(x)$ 与 $g(x)$ 有公根 (在复数域中) 的充分必要条件是它们的结式 $R(f,g)=0$.

\noindent\textbf{推论 5.10.1}\quad 多项式 $f(x)$ 与 $g(x)$ 互素的充分必要条件是 $R(f,g)\neq 0$.

多项式的结式也可以用它们的根来表示.

\noindent\textbf{定理 5.10.2}\quad 设
\[
f(x)=a_0x^n+a_1x^{n-1}+\cdots+a_n,
\]
\[
g(x)=b_0x^m+b_1x^{m-1}+\cdots+b_m,
\]
$f(x)$ 的根为 $x_1,x_2,\cdots,x_n$，$g(x)$ 的根为 $y_1,y_2,\cdots,y_m$，则 $f(x)$ 与 $g(x)$ 的结式为
\begin{equation}
R(f,g)=a_0^mb_0^n\prod_{j=1}^m\prod_{i=1}^n(x_i-y_j).
\tag{5.10.2}
\end{equation}

为了证明上述定理，我们先证明一个引理.

\noindent\textbf{引理 5.10.2}\quad 记号同定理 5.10.2，设 $\lambda$ 是任意的数，则
\[
R(f(x),g(x)(x-\lambda))=(-1)^nf(\lambda)R(f,g),\quad R(f(x),x-\lambda)=(-1)^nf(\lambda).
\]

\noindent\textbf{证明}\quad 由假设，
\[
g(x)(x-\lambda)=b_0x^{m+1}+(b_1-b_0\lambda)x^m+\cdots+(b_m-b_{m-1}\lambda)x-b_m\lambda.
\]
为了书写方便，引入以下记号：
\[
f_0(x)=a_0,\quad f_1(x)=a_0x+a_1,\quad f_2(x)=a_0x^2+a_1x+a_2,\cdots,\quad f_n(x)=f(x).
\]
由结式的定义，
\[
R(f(x),g(x)(x-\lambda))=
\begin{vmatrix}
a_0&a_1&a_2&\cdots&\cdots&a_n&0&\cdots&0\\
0&a_0&a_1&\cdots&\cdots&a_{n-1}&a_n&\cdots&0\\
0&0&a_0&\cdots&\cdots&a_{n-2}&a_{n-1}&\cdots&0\\
\vdots&\vdots&\vdots&\vdots&\vdots&\vdots&\vdots&\vdots&\vdots\\
0&0&\cdots&0&a_0&\cdots&\cdots&\cdots&a_n\\
b_0&b_1-b_0\lambda&b_2-b_1\lambda&\cdots&\cdots&\cdots&-b_m\lambda&\cdots&0\\
0&b_0&b_1-b_0\lambda&\cdots&\cdots&\cdots&b_m-b_{m-1}\lambda&-b_m\lambda&\cdots\\
\vdots&\vdots&\vdots&\vdots&\vdots&\vdots&\vdots&\vdots&\vdots\\
0&\cdots&0&b_0&b_1-b_0\lambda&\cdots&\cdots&\cdots&-b_m\lambda
\end{vmatrix}.
\]
将第 1 列乘以 $\lambda$ 加到第 2 列；再将第 2 列乘以 $\lambda$ 加到第 3 列；$\cdots\cdots$；最后将第 $n+m$ 列乘以 $\lambda$ 加到第 $n+m+1$ 列，则
\[
\begin{aligned}
\text{上式}=
\begin{vmatrix}
f_0(\lambda)&f_1(\lambda)&f_2(\lambda)&\cdots&\cdots&f(\lambda)&\lambda f(\lambda)&\cdots&\lambda^mf(\lambda)\\
0&f_0(\lambda)&f_1(\lambda)&\cdots&\cdots&f_{n-1}(\lambda)&f(\lambda)&\cdots&\lambda^{m-1}f(\lambda)\\
0&0&f_0(\lambda)&\cdots&\cdots&f_{n-2}(\lambda)&f_{n-1}(\lambda)&\cdots&\lambda^{m-2}f(\lambda)\\
\vdots&\vdots&\vdots&\vdots&\vdots&\vdots&\vdots&\vdots&\vdots\\
0&0&\cdots&0&f_0(\lambda)&\cdots&\cdots&\cdots&f(\lambda)\\
b_0&b_1&b_2&\cdots&\cdots&b_m&0&\cdots&0\\
0&b_0&b_1&\cdots&\cdots&b_{m-1}&b_m&0&\cdots\\
\vdots&\vdots&\vdots&\vdots&\vdots&\vdots&\vdots&\vdots&\vdots\\
0&\cdots&0&b_0&b_1&\cdots&\cdots&b_m&0
\end{vmatrix}.
\end{aligned}
\]
将第 2 行乘以 $-\lambda$ 加到第 1 行；再将第 3 行乘以 $-\lambda$ 加到第 2 行；$\cdots\cdots$；最后将第 $m+1$ 行乘以 $-\lambda$ 加到第 $m$ 行，则
\[
\text{上式}=
\begin{vmatrix}
a_0&a_1&a_2&\cdots&\cdots&a_n&0&\cdots&0\\
0&a_0&a_1&\cdots&\cdots&a_{n-1}&a_n&\cdots&0\\
\vdots&\vdots&\vdots&\vdots&\vdots&\vdots&\vdots&\vdots&\vdots\\
0&\cdots&0&a_0&a_1&\cdots&\cdots&a_n&0\\
0&0&\cdots&0&f_0(\lambda)&\cdots&\cdots&\cdots&f(\lambda)\\
b_0&b_1&b_2&\cdots&\cdots&b_m&0&\cdots&0\\
0&b_0&b_1&\cdots&\cdots&b_{m-1}&b_m&0&\cdots\\
\vdots&\vdots&\vdots&\vdots&\vdots&\vdots&\vdots&\vdots&\vdots\\
0&\cdots&0&b_0&b_1&\cdots&\cdots&b_m&0
\end{vmatrix}.
\]
将行列式按第 $n+m+1$ 列展开，则
\[
\text{上式}=(-1)^nf(\lambda)
\begin{vmatrix}
a_0&a_1&a_2&\cdots&\cdots&a_n&0&\cdots\\
0&a_0&a_1&\cdots&\cdots&a_{n-1}&a_n&\cdots\\
\vdots&\vdots&\vdots&\vdots&\vdots&\vdots&\vdots&\vdots\\
0&\cdots&0&a_0&a_1&\cdots&\cdots&a_n\\
b_0&b_1&b_2&\cdots&\cdots&b_m&0&\cdots\\
0&b_0&b_1&\cdots&\cdots&b_{m-1}&b_m&\cdots\\
\vdots&\vdots&\vdots&\vdots&\vdots&\vdots&\vdots&\vdots\\
0&\cdots&0&b_0&b_1&\cdots&\cdots&b_m
\end{vmatrix}
=(-1)^nf(\lambda)R(f,g).
\]
因此 $R(f(x),g(x)(x-\lambda))=(-1)^nf(\lambda)R(f,g)$. 同理可证 $R(f(x),x-\lambda)=(-1)^nf(\lambda)$. $\square$

\noindent\textbf{定理 5.10.2 的证明}\quad 令 $f(x)=a_0f_1(x),\ g(x)=b_0g_1(x)$，则 $f_1(x)$ 与 $f(x)$ 有相同的根，$g_1(x)$ 与 $g(x)$ 有相同的根. 由结式的定义知道
\[
R(f,g)=a_0^mb_0^nR(f_1,g_1).
\]
反复利用引理 5.10.2 可得
\[
\begin{aligned}
R(f_1,g_1)
&=R(f_1,(x-y_1)(x-y_2)\cdots(x-y_m))\\
&=(-1)^nf_1(y_1)R(f_1,(x-y_2)\cdots(x-y_m))=\cdots\\
&=\prod_{j=1}^m\left((-1)^nf_1(y_j)\right)
=\prod_{j=1}^m\left((-1)^n\prod_{i=1}^n(y_j-x_i)\right)\\
&=\prod_{j=1}^m\prod_{i=1}^n(x_i-y_j).
\end{aligned}
\]
由此即得 (5.10.2) 式. $\square$

利用结式，可定义一个多项式的判别式如下.

\noindent\textbf{定义 5.10.2}\quad 多项式
\[
f(x)=a_0x^n+a_1x^{n-1}+\cdots+a_n
\]
的判别式定义为
\[
\Delta(f)=(-1)^{\frac12n(n-1)}a_0^{-1}R(f,f').
\]

\noindent\textbf{定理 5.10.3}\quad 多项式
\[
f(x)=a_0x^n+a_1x^{n-1}+\cdots+a_n
\]
的判别式
\begin{equation}
\Delta(f)=a_0^{2n-2}\prod_{1\leq i<j\leq n}(x_i-x_j)^2,
\tag{5.10.3}
\end{equation}
其中 $x_1,x_2,\cdots,x_n$ 为 $f(x)$ 的根.

\noindent\textbf{证明}\quad 由 (5.10.2) 式知道
\[
R(f,g)=a_0^m\prod_{i=1}^n g(x_i).
\]
现令 $g(x)=f'(x),\ \deg f'(x)=n-1$，因此
\[
R(f,f')=a_0^{n-1}\prod_{i=1}^n f'(x_i).
\]
又
\[
f(x)=a_0(x-x_1)(x-x_2)\cdots(x-x_n),
\]
\[
f'(x)=a_0\left(\sum_{j=1}^n(x-x_1)\cdots(x-x_{j-1})(x-x_{j+1})\cdots(x-x_n)\right),
\]
\[
f'(x_i)=a_0(x_i-x_1)\cdots(x_i-x_{i-1})(x_i-x_{i+1})\cdots(x_i-x_n).
\]
因此
\[
R(f,f')=(-1)^{\frac12n(n-1)}a_0^{2n-1}\prod_{1\leq i<j\leq n}(x_i-x_j)^2.
\]
由此即得 (5.10.3) 式. $\square$

例如，$ax^2+bx+c$ 的判别式为 $a^2(x_1-x_2)^2=b^2-4ac$. 这是我们早已熟悉的判别式.

\noindent\textbf{推论 5.10.2}\quad 多项式 $f(x)$ 有重根的充分必要条件是它的判别式 $\Delta(f)=0$.

作为结式的应用，我们来求解二元高次方程组. 设
\begin{equation}
\begin{cases}
f(x,y)=0,\\
g(x,y)=0
\end{cases}
\tag{5.10.4}
\end{equation}
是由两个二元多项式组成的方程组. 我们的目的是把求解这组方程先归结为求解一个一元高次方程. 将 $f(x,y),g(x,y)$ 整理为关于 $x$ 的多项式：
\begin{equation}
f(x,y)=a_0(y)x^n+a_1(y)x^{n-1}+\cdots+a_n(y),
\tag{5.10.5}
\end{equation}
\begin{equation}
g(x,y)=b_0(y)x^m+b_1(y)x^{m-1}+\cdots+b_m(y),
\tag{5.10.6}
\end{equation}
其中 $a_i(y),b_j(y)$ 都是 $y$ 的多项式且 $a_0(y)\neq 0,\ b_0(y)\neq 0$. 令
\[
R_x(f,g)=
\begin{vmatrix}
a_0(y)&a_1(y)&a_2(y)&\cdots&\cdots&a_n(y)&0&\cdots&0\\
0&a_0(y)&a_1(y)&\cdots&\cdots&a_{n-1}(y)&a_n(y)&\cdots&0\\
0&0&a_0(y)&\cdots&\cdots&a_{n-2}(y)&a_{n-1}(y)&\cdots&0\\
\vdots&\vdots&\vdots&\vdots&\vdots&\vdots&\vdots&\vdots&\vdots\\
0&0&\cdots&0&a_0(y)&\cdots&\cdots&\cdots&a_n(y)\\
b_0(y)&b_1(y)&b_2(y)&\cdots&\cdots&\cdots&b_m(y)&\cdots&0\\
0&b_0(y)&b_1(y)&\cdots&\cdots&\cdots&b_{m-1}(y)&b_m(y)&\cdots\\
\vdots&\vdots&\vdots&\vdots&\vdots&\vdots&\vdots&\vdots&\vdots\\
0&\cdots&0&b_0(y)&b_1(y)&\cdots&\cdots&\cdots&b_m(y)
\end{vmatrix}.
\]
这是一个关于 $y$ 的多项式. 如果 $(\alpha,\beta)$ 是方程组 (5.10.4) 的解，则 $\alpha$ 是 $f(x,\beta),g(x,\beta)$ 的公根，从而 $\beta$ 是 $R_x(f,g)$ 的根. 因此，我们可先求出 $R_x(f,g)=0$ 的所有根 $\beta_i$，再代入 (5.10.5) 式、(5.10.6) 式. 这时，或 $a_0(\beta_i)=b_0(\beta_i)=0$，或存在 $\alpha_i$，使 $(\alpha_i,\beta_i)$ 是方程组 (5.10.4) 的解. 由此即可求出方程组 (5.10.4) 的一切解. 对于多于两个未知数的高次方程组，也可采用类似方法逐个“消去”未知数，从而求出方程组的解.

\subsection*{习题 5.10}

1. 计算下列多项式的结式：

(1) $f(x)=x^3+3x^2-x+4,\ g(x)=x^2-2x-1$；

(2) $f(x)=x^5+1,\ g(x)=x^2+x+1$.

2. 计算 $x^3+px+q$ 的判别式.

3. 设 $f(x)=a_0x^n+a_1x^{n-1}+\cdots+a_n,\ g(x)=b_0x^m+b_1x^{m-1}+\cdots+b_m$. 求证：

(1) $R(f,g)=(-1)^{mn}R(g,f)$；

(2) 若 $a,b$ 为常数，则 $R(af,bg)=a^mb^nR(f,g)$.

4. 求证：$R(f,g_1g_2)=R(f,g_1)R(f,g_2)$.

5. 设 $f(x)$ 是实系数多项式，求证：

(1) 若 $\Delta(f)<0$，则 $f(x)$ 无重根且有奇数对虚根；

(2) 若 $\Delta(f)>0$，则 $f(x)$ 无重根且有偶数对虚根.

6. 设 $f(x)$ 是数域 $\mathbb{K}$ 上的多项式，已知 $\Delta(f(x))$，试求 $\Delta(f(x^2))$.

7. 求下列参数曲线的直角坐标方程：
\[
\begin{cases}
x=t^3+2t-3,\\
y=t^2-t+1.
\end{cases}
\]

\section*{历史与展望}

一元多项式 (一元 $n$ 次方程) 是人类最早开始研究的数学对象之一. 4000 年前，古巴比伦人用楔形文字记录了他们求解包含一个未定元的某些线性方程和二次方程. 公元前 6 世纪，古希腊人用几何的方法求解了类似的方程. 例如，Euclid 的著作《几何原本》第 6 卷的命题 28 和命题 29 实际上给出了二次方程 $x(a-x)=S$ 的求解方法.

公元 3 世纪，Diophantus (丢番图) 将研究范围扩大到很多其他类型的方程，包括高次方程、多变元方程以及同类方程的方程组. 他还发展了第一个字母符号体系以记述这些方程和相关的代数问题. 公元 8 世纪，阿拉伯学者开始把方程作为有价值的研究对象，同时根据在已有技术下求解方程的难易程度，对线性方程、二次方程和三次方程进行了分类. Khwarizmi (花拉子密) 在其著作中首次使用了“代数 (Algebra)”这个词语.

三次方程和四次方程解法的发现是发生在意大利的一段充满着商业与竞争意味的历史. Pacioli (帕乔利) 曾在其著作《数学大全》中列出两类可能没有解的三次方程：
\[
(1)\ n=ax+bx^3,\qquad (2)\ n=ax^2+bx^3.
\]
16 世纪初，Ferro (费罗) 发现了上面第一类三次方程的解法，他在去世之前把这个解法告诉了他的学生 Fiore (菲奥雷). 1535 年 Tartaglia (塔尔塔利亚) 得到了上面两类三次方程的解法，并在与 Fiore 的数学竞赛中获得了胜利. Cardano 在得知这一消息后，开始了与 Tartaglia 长达三个月的书信交往，邀请他来家中做客并热情款待. 1539 年 Tartaglia 用 25 行诗的形式写下了第一类三次方程的解法，并让 Cardano 发誓绝不泄露这个秘密. Cardano 在得到这个秘密后，进一步研究得到了三次方程的一般解法. 在这一探索过程中，他的助手 Ferrari 于 1540 年发现了四次方程的一般解法. 1545 年 Cardano 出版了著作《大术》，包含了三次方程和四次方程的一般解法，《大术》也是第一个严谨地提出了复数概念的数学文献.

法国数学家 Descartes 在 1637 年出版的著作《几何学》中，在平面上引入了坐标系，这一思想促使几何学代数化，建立了解析几何这一数学分支. Descartes 的另外一个重要贡献是他给我们带来了现代字母符号体系，用字母表开头的几个小写字母表示已知数，字母表末尾的几个小写字母表示未知数，从此方程 (多项式) 有了简洁方便的现代写法.

法国数学家 Vieta 是方程研究的先驱. 在 1615 年 (他去世后 12 年) 发表的《论方程的整理和修正》的第二篇论文中，Vieta 给出了一元五次以下方程的根与系数之间关系的公式 (称为 Vieta 定理)，后来 Girard (吉拉德) 在《代数新发现》一书中把这些公式推广到任意次方程. 英
国数学家 Newton 在 1665 年或 1666 年写下了一些简短的笔记 (收录在他的《数学全集》第一卷)，其中隐含了对称多项式基本定理 (也称为 Newton 定理). 由这两个定理即可推出：一元 $n$ 次方程的解的对称多项式可以表示为方程系数的有理函数. 这一研究方程解的对称性的思想，为解决五次及五次以上方程的根式求解问题开辟了道路.

瑞士数学家 Euler (欧拉) 在 1732 年首次研究了一般五次方程的求解问题. 在论文《论任意次方程的解》中，Euler 提出任意 $n$ 次方程解的表达式也许是这样的形式：
\[
A+B\sqrt[n]{\alpha}+C(\sqrt[n]{\alpha})^2+D(\sqrt[n]{\alpha})^3+\cdots,
\]
其中 $\alpha$ 是某个 $n-1$ 次“辅助方程”的解，而 $A,B,C,D,\cdots$ 是方程系数的有理函数. 法国数学家 Vandermonde 提出了一个非常重要的见解：把方程的解写成所有解的对称多项式或部分对称多项式的表达式，这里部分对称多项式是指在所有解的部分置换下保持不变的多项式.

法国数学家 Lagrange 在 1771 年的论文《对方程的代数解的思考》中，把通过方程解的置换来求解方程的想法呈现给世人，这与 Vandermonde 的想法是类似的，但更加透彻和深刻. 以三次方程为例，设 $\alpha,\beta,\gamma$ 是缺项三次方程 $x^3+px+q=0$ 的三个解，$\omega=\cos\dfrac{2\pi}{3}+\mathrm{i}\sin\dfrac{2\pi}{3}$ 是 3 次单位根. 构造解的表达式 $\alpha+\omega\beta+\omega^2\gamma$ (称为 Lagrange 预解式)，当置换解 $\alpha,\beta,\gamma$ 时，上述表达式取六个不同的值，分别为
\[
t_1=\alpha+\omega\beta+\omega^2\gamma,\quad
t_2=\alpha+\omega^2\beta+\omega\gamma,\quad
t_3=\omega^2t_2,\quad
t_4=\omega t_2,\quad
t_5=\omega^2t_1,\quad
t_6=\omega t_1.
\]
构造以 $t_i$ 为根的六次方程 (称为 Lagrange 预解方程)：
\[
(x-t_1)(x-t_2)(x-t_3)(x-t_4)(x-t_5)(x-t_6)=0,
\]
化简可得 $(x^3-t_1^3)(x^3-t_2^3)=0$. 这实际上是关于 $t_1^3,t_2^3$ 的二次方程，求出 $t_1,t_2$ 之后便可得到三次方程的解 $\alpha,\beta,\gamma$ 的表达式，即 Cardano 公式. Lagrange 对四次方程也进行了同样的处理，他得到了一个 24 次预解方程，降次之后成为一个六次方程，实际上这是一个关于 $x^2$ 的三次方程，从而是可解的. 这个过程完全复原了四次方程的 Ferrari 解法. 五个对象的置换一共有 $5!=120$ 个，因此五次方程的预解方程是 120 次方程，Lagrange 通过一些技巧把这个预解方程降次成为一个 24 次方程，然后他便陷入了困境. 即使如此，Lagrange 还是抓住了方程可解性的重点：必须深入研究方程解的置换，以及预解方程在置换作用下发生的变化. Lagrange 还证明了一个重要的定理，也就是现今群论中的 Lagrange 定理，它的原始表达为：对一个 $n$ 元多项式的 $n$ 个未定元进行置换，得到的互异多项式的个数一定要整除置换的总个数 $n!$.

意大利数学家 Ruffini (鲁菲尼) 遵循 Lagrange 的想法，仔细观察了置换未定元时互异多项式的可能个数，然后证明了对于一般的五次方程，不可能得到一个三次或四次预解方程，再由 Lagrange 的结果便可得到一般的五次方程不能根式求解，即一般的五次方程的解不能写成方程系数的加、减、乘、除和开方的表达式. Ruffini 的第一个证明存在缺陷，但他继续研究，至少发表了 3 个证明. 他把这些证明寄给一些著名数学家 (包括 Lagrange)，但不是被忽视，就是被拒绝接受. 直到 1821 年，Ruffini 收到了法国数学家 Cauchy 的来信，Cauchy 在信中赞扬了 Ruffini 的工作，并认为他已经证明了一般的五次方程不能根式求解. 尽管如此，同时代的数学家还是认为 Ruffini 的证明有缺陷，且他书写证明的风格令人难以理解，这正是 Lagrange 否认他的原因.

1820 年，年仅 18 岁的挪威数学家 Abel (阿贝尔) 开始着手研究一般的五次方程，并且证明了不能根式求解的定理. 1824 年，他自己出钱印刷发表了这个证明. Abel 的证明结合了他从 Euler、Lagrange、Ruffini 和 Cauchy 等人那里学到的想法，并用其独创的方法和见解将这些想法结合在一起. 他用的是反证法：从假设欲证的结论不成立开始，然后论证这将导致逻辑矛盾. Abel 的证明经历了很长的时间才广为人知，但他的证明并没有终结一元 $n$ 次方程的求解理论. 虽然一般的五次方程不能根式求解，但一些特殊的五次方程可以根式求解. 因此可以自然地问：满足怎样条件的五次方程可以根式求解呢？五次以上方程的情况又如何呢？给出这些问题完整解答的是法国数学家 Galois.

一个一元 $n$ 次方程的系数属于某个域，但它在这个域中可能没有解. 为了包含它的那些解，我们把系数域扩张为一个更大的域，称为分裂域. 1829 年，年仅 18 岁的 Galois 发现方程解的形式取决于系数域与分裂域之间的关系，这种关系可以用群论的语言来表达. 简单来说，方程解之间的置换诱导了分裂域的自同构，这些自同构保持系数域不变，它们构成了一个群，称为方程的 Galois 群. Galois 证明了：方程可以用根式求解当且仅当方程的 Galois 群是可解群；一般的五次及五次以上方程的 Galois 群不是可解群，从而不能用根式求解. 1830--1831 年，Galois 关于方程解的论文的前三版分别提交给了法国科学院的 Cauchy、Fourier (傅里叶) 和 Poisson (泊松)，不过都被忽略或拒稿. 1832 年 5 月 30 日，Galois 因决斗而亡. 10 年之后，法国数学家 Liouville (刘维尔) 对 Galois 的论文产生了兴趣. 1843 年他在法国科学院宣读了 Galois 的主要研究成果，1846 年他又在自己创办的数学杂志上发表了 Galois 的所有论文. 直到此时，Galois 理论才被数学界广为所知. Galois 理论不仅完美解决了一元 $n$ 次方程的根式求解问题，而且它将“群、域”等抽象的新数学对象引入了代数学的研究视野，拉开了现代代数学研究的序幕.

\section*{复习题五}

1. 设 $f(x),g(x)$ 是次数不小于 1 的互素多项式，求证：必唯一地存在两个多项式 $u(x),v(x)$，使
\[
f(x)u(x)+g(x)v(x)=1,
\]
且 $\deg u(x)<\deg g(x),\ \deg v(x)<\deg f(x)$.

2. 求证：$(f(x),g(x))=1$ 的充分必要条件是对任意的正整数 $m,n$，$(f(x)^m,g(x)^n)=1$.

3. 设 $f(x)=x^m-1,\ g(x)=x^n-1$，求证：$(f(x),g(x))=x^d-1$，其中 $d$ 是 $m,n$ 的最大公因子.

4. 设 $(f(x),g(x))=d(x)$，求证：对任意的正整数 $n$，
\[
(f(x)^n,f(x)^{n-1}g(x),\cdots,f(x)g(x)^{n-1},g(x)^n)=d(x)^n.
\]

5. 设 $f(x)$ 是数域 $\mathbb{K}$ 上的非常数多项式，求证：$f(x)$ 等于某个不可约多项式的幂的充分必要条件是对任意的非常数多项式 $g(x)$，或者 $f(x)$ 和 $g(x)$ 互素，或者 $f(x)$ 可以整除 $g(x)$ 的某个幂.

6. 设 $p(x)$ 是数域 $\mathbb{K}$ 上的不可约多项式，$f(x)$ 是 $\mathbb{K}$ 上的多项式. 证明：若 $p(x)$ 的某个复根 $\alpha$ 也是 $f(x)$ 的根，则 $p(x)\mid f(x)$. 特别，$p(x)$ 的任一复根都是 $f(x)$ 的根.

7. 设 $f(x)$ 是非常数多项式且 $f(x)$ 可以整除 $f(x^m)(m>1)$，求证：$f(x)$ 的根只能是 0 或 1 的某个方根.

8. 求证：$f(x)=\sin x$ 在实数域内不能表示为 $x$ 的多项式.

9. 设 $n$ 是奇数，求证：$(x+y)(y+z)(x+z)$ 可整除 $(x+y+z)^n-x^n-y^n-z^n$.

10. 设 $x_1,x_2,x_3$ 是三次方程 $x^3+px^2+qx+r=0\ (r\neq 0)$ 的 3 个根，求这 3 个根倒数的平方和.

11. 设 $f(x)=a_nx^n+a_{n-1}x^{n-1}+\cdots+a_1x+a_0\ (a_na_0\neq 0)$ 是数域 $\mathbb{K}$ 上的可约多项式，求证：多项式 $g(x)=a_0x^n+a_1x^{n-1}+\cdots+a_{n-1}x+a_n$ 在 $\mathbb{K}$ 上也可约.

12. 求证：实系数方程 $x^3+px^2+qx+r=0$ 的根的实部全是负数的充分必要条件是
\[
p>0,\quad r>0,\quad pq>r.
\]

13. 设 $\varepsilon=\cos\dfrac{2\pi}{n}+\mathrm{i}\sin\dfrac{2\pi}{n}$ 是 1 的 $n$ 次根，求证：$\varepsilon^{mi}(i=1,2,\cdots,n)$ 是 $x^n-1=0$ 的全部根的充分必要条件是 $(m,n)=1$.

14. 设 $f(x)$ 是实系数首一多项式且无实数根，求证：$f(x)$ 可以表示为两个实系数多项式的平方和.

15. 设 $f(x)$ 是实系数多项式，若对任意的有理数 $c,\ f(c)$ 总是有理数，求证：$f(x)$ 是有理系数多项式.

16. 设 $f(x)$ 是有理系数多项式，$a,b,c$ 是有理数，但 $\sqrt c$ 是无理数. 求证：若 $a+b\sqrt c$ 是 $f(x)$ 的根，则 $a-b\sqrt c$ 也是 $f(x)$ 的根.

17. 设 $f(x)$ 是有理系数多项式，$a,b,c,d$ 是有理数，但 $\sqrt c,\sqrt d,\sqrt{cd}$ 都是无理数. 求证：若 $a\sqrt c+b\sqrt d$ 是 $f(x)$ 的根，则下列数也是 $f(x)$ 的根：
\[
a\sqrt c-b\sqrt d,\quad -a\sqrt c+b\sqrt d,\quad -a\sqrt c-b\sqrt d.
\]

18. 求以 $\sqrt2+\sqrt[3]{3}$ 为根的次数最小的首一有理系数多项式.

19. 求证：有理系数多项式 $x^4+px^2+q$ 在有理数域上可约的充分必要条件是或者 $p^2-4q=k^2$，其中 $k$ 是一个有理数；或者 $q$ 是某个有理数的平方，且 $\pm2\sqrt q-p$ 也是有理数的平方.

20. 设 $f(x)$ 是有理系数多项式，已知 $\sqrt[n]{2}$ 是 $f(x)$ 的根，证明：$\sqrt[n]{2}\varepsilon,\sqrt[n]{2}\varepsilon^2,\cdots,\sqrt[n]{2}\varepsilon^{n-1}$ 也是 $f(x)$ 的根，其中 $\varepsilon=\cos\dfrac{2\pi}{n}+\mathrm{i}\sin\dfrac{2\pi}{n}$ 是 1 的 $n$ 次根.

21. 设 $f(x)$ 是次数大于 1 的奇数次有理系数不可约多项式，求证：若 $x_1,x_2$ 是 $f(x)$ 在复数域内的两个不同的根，则 $x_1+x_2$ 必不是有理数.

22. 设 $f(x)=(x-a_1)(x-a_2)\cdots(x-a_n)-1$，其中 $a_1,a_2,\cdots,a_n$ 是 $n$ 个不同的整数，求证：$f(x)$ 在有理数域上不可约.

23. 设 $f(x)=(x-a_1)^2(x-a_2)^2\cdots(x-a_n)^2+1$，其中 $a_1,a_2,\cdots,a_n$ 是 $n$ 个不同的整数，求证：$f(x)$ 在有理数域上不可约.

24. 求 $n$ 次多项式 $x^n+px+q(n>1)$ 的判别式.

25. 求下列多项式的判别式：

(1) $f(x)=x^n+2x+1$；\qquad (2) $f(x)=x^n+2$；

(3) $f(x)=x^{n-1}+x^{n-2}+\cdots+x+1$.

26. 求证：多项式 $x^4+px+q$ 有重因式的充分必要条件是 $27p^4=256q^3$.

27. 设 $f(x)$ 和 $g(x)$ 是次数大于 1 的多项式，求证：
\[
\Delta(f(x)g(x))=\Delta(f(x))\Delta(g(x))R(f,g)^2.
\]

28. 设 $g(x)$ 是次数大于 1 的多项式，求证：
\[
\Delta((x-a)g(x))=g(a)^2\Delta(g(x)).
\]

29. 设 $f(x)=g(h(x))$，其中 $h(x)$ 是 $m$ 次首一多项式，$g(x)$ 是 $n$ 次首一多项式，其根为 $x_1,x_2,\cdots,x_n$，求证：
\[
\Delta(f(x))=\Delta(g(x))^m\Delta(h(x)-x_1)\Delta(h(x)-x_2)\cdots\Delta(h(x)-x_n).
\]

30. 求下列参数曲线的直角坐标方程：
\[
\begin{cases}
x=2(t+1)/(t^2+1),\\
y=t^2/(2t-1).
\end{cases}
\]

31. 设 $f(x),g(x)$ 是数域 $\mathbb{K}$ 上的互素多项式，$A$ 是 $\mathbb{K}$ 上的 $n$ 阶方阵，满足 $f(A)=O$，证明：$g(A)$ 是可逆阵.

32. 设 $f(x),g(x)$ 是数域 $\mathbb{K}$ 上的互素多项式，$A$ 是 $\mathbb{K}$ 上的 $n$ 阶方阵，证明：$f(A)g(A)=O$ 的充分必要条件是 $r(f(A))+r(g(A))=n$.

33. 设 $f(x),g(x)$ 是数域 $\mathbb{K}$ 上的互素多项式，$\varphi$ 是 $\mathbb{K}$ 上 $n$ 维线性空间 $V$ 上的线性变换，满足 $f(\varphi)g(\varphi)=0$，证明：$V=V_1\oplus V_2$，其中 $V_1=\ker f(\varphi),V_2=\ker g(\varphi)$.

34. 设 $\mathbb{Q}(\sqrt[n]{2})=\{a_0+a_1\sqrt[n]{2}+a_2\sqrt[n]{4}+\cdots+a_{n-1}\sqrt[n]{2^{n-1}}\mid a_i\in\mathbb{Q},0\leq i\leq n-1\}$，证明：$\mathbb{Q}(\sqrt[n]{2})$ 是一个数域，并求 $\mathbb{Q}(\sqrt[n]{2})$ 作为 $\mathbb{Q}$ 上线性空间的一组基.

35. 设 $f(x)=x^n+a_1x^{n-1}+\cdots+a_{n-1}x+a_n$ 是数域 $\mathbb{K}$ 上的不可约多项式，$\varphi$ 是 $\mathbb{K}$ 上 $n$ 维线性空间 $V$ 上的线性变换，$\alpha_1\neq 0,\alpha_2,\cdots,\alpha_n$ 是 $V$ 中的向量，满足：
\[
\varphi(\alpha_1)=\alpha_2,\ \varphi(\alpha_2)=\alpha_3,\ \cdots,\ \varphi(\alpha_{n-1})=\alpha_n,\ \varphi(\alpha_n)=-a_n\alpha_1-a_{n-1}\alpha_2-\cdots-a_1\alpha_n.
\]
证明：$\{\alpha_1,\alpha_2,\cdots,\alpha_n\}$ 是 $V$ 的一组基.

\end{document}
